% ============================================================================
% Aleph Filter: Understanding and Implementation
% A Technical Report by Neel
% ============================================================================
\documentclass[11pt, a4paper]{report}

% === GEOMETRY & LAYOUT ===
\usepackage[
  top=1in,
  bottom=1in,
  left=1.15in,
  right=1.15in
]{geometry}

% === FONTS & ENCODING ===
\usepackage[T1]{fontenc}
\usepackage[expansion=false]{microtype}  % Better text justification

% === MATH ===
\usepackage{amsmath, amssymb, amsthm}

% === COLORS ===
\usepackage[dvipsnames, svgnames]{xcolor}

% Custom color palette
\definecolor{AccentBlue}{HTML}{2563EB}
\definecolor{AccentTeal}{HTML}{0D9488}
\definecolor{AccentOrange}{HTML}{EA580C}
\definecolor{AccentRed}{HTML}{DC2626}
\definecolor{AccentGreen}{HTML}{16A34A}
\definecolor{LightGray}{HTML}{F8FAFC}
\definecolor{MedGray}{HTML}{94A3B8}
\definecolor{DarkGray}{HTML}{334155}
\definecolor{CodeBg}{HTML}{F1F5F9}
\definecolor{BitZero}{HTML}{E2E8F0}
\definecolor{BitOne}{HTML}{BFDBFE}
\definecolor{VoidColor}{HTML}{FED7AA}
\definecolor{TombstoneColor}{HTML}{FECACA}

% === HEADERS & FOOTERS ===
\usepackage{fancyhdr}
\pagestyle{fancy}
\fancyhf{}
\fancyhead[L]{\small\textcolor{black}{\leftmark}}
\fancyhead[R]{\small\textcolor{black}{Aleph Filter}}
\fancyfoot[C]{\small\textcolor{black}{\thepage}}
\renewcommand{\headrulewidth}{0.4pt}
\renewcommand{\footrulewidth}{0pt}

% === CHAPTER & SECTION STYLING ===
\usepackage{titlesec}

\titleformat{\chapter}[display]
  {\normalfont\huge\bfseries\color{black}}
  {\chaptertitlename\ \thechapter}{18pt}{\Huge}
\titlespacing*{\chapter}{0pt}{-20pt}{30pt}

\titleformat{\section}
  {\normalfont\Large\bfseries\color{black}}
  {\thesection}{1em}{}

\titleformat{\subsection}
  {\normalfont\large\bfseries\color{black}}
  {\thesubsection}{1em}{}

% === HYPERLINKS ===
\usepackage[
  colorlinks=true,
  linkcolor=black,
  citecolor=black,
  urlcolor=black
]{hyperref}
\usepackage{xurl}

% === FIGURES & CAPTIONS ===
\usepackage{graphicx}
\usepackage{float}
\usepackage[
  font=small,
  labelfont=bf,
  labelsep=period,
  justification=centering
]{caption}

% === LISTS ===
\usepackage{enumitem}
\setlist[itemize]{topsep=4pt, itemsep=2pt}
\setlist[enumerate]{topsep=4pt, itemsep=2pt}

% === CODE LISTINGS ===
\usepackage{listings}

\lstdefinestyle{rust}{
  language=C,  % Close enough base for Rust
  basicstyle=\small\ttfamily,
  backgroundcolor=\color{CodeBg},
  keywordstyle=\color{AccentBlue}\bfseries,
  stringstyle=\color{AccentGreen},
  commentstyle=\color{MedGray}\itshape,
  numberstyle=\tiny\color{MedGray},
  numbers=left,
  numbersep=8pt,
  breaklines=true,
  frame=single,
  framerule=0pt,
  framexleftmargin=6pt,
  xleftmargin=14pt,
  tabsize=4,
  showstringspaces=false,
  morekeywords={fn, let, mut, pub, struct, impl, self, Self, use, mod, 
                 crate, u8, u16, u32, u64, usize, bool, true, false,
                 Option, Some, None, Vec, String, Result, Ok, Err,
                 if, else, for, while, loop, match, return, as, const},
  literate={->}{$\rightarrow$}{2}
           {=>}{$\Rightarrow$}{2},
  belowskip=8pt,
  aboveskip=8pt,
}
\lstset{style=rust}

% === ALGORITHMS / PSEUDOCODE ===
\usepackage[
  ruled,
  vlined,
  linesnumbered
]{algorithm2e}
\SetAlgoLined
\SetKwComment{Comment}{// }{}

% === BOXES & CALLOUTS ===
\usepackage[most]{tcolorbox}

% Definition box
\newtcolorbox{definitionbox}[1][]{
  colback=blue!3,
  colframe=AccentBlue,
  coltitle=white,
  fonttitle=\bfseries,
  title={#1},
  arc=2pt,
  boxrule=1pt,
  left=8pt, right=8pt, top=6pt, bottom=6pt,
}

% Key insight box
\newtcolorbox{insightbox}[1][Key Insight]{
  colback=teal!3,
  colframe=AccentTeal,
  coltitle=white,
  fonttitle=\bfseries,
  title={#1},
  arc=2pt,
  boxrule=1pt,
  left=8pt, right=8pt, top=6pt, bottom=6pt,
}

% Warning / gotcha box
\newtcolorbox{warningbox}[1][Watch Out]{
  colback=orange!3,
  colframe=AccentOrange,
  coltitle=white,
  fonttitle=\bfseries,
  title={#1},
  arc=2pt,
  boxrule=1pt,
  left=8pt, right=8pt, top=6pt, bottom=6pt,
}

% Example box
\newtcolorbox{examplebox}[1][Example]{
  colback=green!2,
  colframe=AccentGreen,
  coltitle=white,
  fonttitle=\bfseries,
  title={#1},
  arc=2pt,
  boxrule=1pt,
  left=8pt, right=8pt, top=6pt, bottom=6pt,
}

% === TIKZ (diagrams, arrays, arrows) ===
\usepackage{tikz}
\usetikzlibrary{
  arrows.meta,       % Better arrowheads
  positioning,       % Relative positioning
  calc,              % Coordinate calculations
  decorations.pathreplacing,  % Braces
  fit,               % Fit nodes around others
  backgrounds,       % Background layers
  patterns,          % Fill patterns
  chains,            % Node chains
  matrix,            % Matrix of nodes
  shapes.geometric,  % Extra shapes
}

% --- CUSTOM TIKZ STYLES ---

% Slot in an array/table
\tikzstyle{slot} = [
  draw, 
  minimum width=1cm, 
  minimum height=0.7cm,
  font=\small\ttfamily,
  inner sep=2pt,
]

% Empty slot
\tikzstyle{emptyslot} = [slot, fill=BitZero]

% Occupied slot
\tikzstyle{fullslot} = [slot, fill=BitOne]

% Void entry slot
\tikzstyle{voidslot} = [slot, fill=VoidColor]

% Tombstone slot
\tikzstyle{tombstoneslot} = [slot, fill=TombstoneColor]

% Bit in a bit array (smaller)
\tikzstyle{bit} = [
  draw,
  minimum width=0.5cm,
  minimum height=0.5cm,
  font=\small\ttfamily,
  inner sep=1pt,
]

% Hash function box
\tikzstyle{hashbox} = [
  draw,
  rounded corners=3pt,
  fill=AccentBlue!10,
  minimum width=1.8cm,
  minimum height=0.7cm,
  font=\small,
  text=AccentBlue,
]

% Label style
\tikzstyle{lbl} = [font=\small, text=DarkGray]

% Arrow styles
\tikzstyle{arrow} = [-{Stealth[length=6pt]}, thick]
\tikzstyle{dualarrow} = [{Stealth[length=6pt]}-{Stealth[length=6pt]}, thick]
\tikzstyle{dasharrow} = [-{Stealth[length=6pt]}, thick, dashed]

% Brace annotation
\tikzstyle{brace} = [
  decorate,
  decoration={brace, amplitude=6pt, raise=4pt},
  thick,
  DarkGray,
]

% --- CUSTOM COMMANDS ---

% Fingerprint notation
\newcommand{\fp}[1]{\texttt{\textcolor{AccentBlue}{#1}}}

% Quotient notation
\newcommand{\quot}[1]{\texttt{\textcolor{AccentTeal}{#1}}}

% Remainder notation
\newcommand{\rem}[1]{\texttt{\textcolor{AccentOrange}{#1}}}

% Void marker
\newcommand{\void}{\textcolor{AccentOrange}{\textsc{void}}}

% Tombstone marker
\newcommand{\tombstone}{\textcolor{AccentRed}{\textsc{tomb}}}

% Big-O notation
\newcommand{\bigO}[1]{\ensuremath{\mathcal{O}(#1)}}

% False Positive Rate
\newcommand{\FPR}{\textsf{FPR}}

% Highlight inline
\newcommand{\hl}[1]{\colorbox{yellow!25}{#1}}

% === THEOREM ENVIRONMENTS ===
\theoremstyle{definition}
\newtheorem{definition}{Definition}[chapter]
\newtheorem{theorem}{Theorem}[chapter]
\newtheorem{lemma}[theorem]{Lemma}

% === TABLE OF CONTENTS DEPTH ===
\setcounter{tocdepth}{2}

% === BIBLIOGRAPHY ===
% If you want to use biblatex later, uncomment below and create a .bib file:
% \usepackage[backend=biber, style=numeric, sorting=none]{biblatex}
% \addbibresource{references.bib}

% ============================================================================
% DOCUMENT BEGINS
% ============================================================================
\begin{document}

% --- TITLE PAGE ---
\begin{titlepage}
  \centering
  \vspace*{3cm}
  
  % Title block
  {\fontsize{36}{40}\selectfont\bfseries Aleph Filter}\\[0.6cm]
  {\Large\color{DarkGray} To Infinity in Constant Time}
  
  \vspace{0.8cm}
  \rule{0.4\textwidth}{0.8pt}
  \vspace{0.8cm}
  
  {\large Understanding, Analysis, and Rust Implementation}
  
  \vspace{2.5cm}
  
  % Author
  {\Large\bfseries Neel Bansal}\\[0.4cm]
  {\color{black} \today}
  
  \vfill
  
  % Bottom box
  \begin{tcolorbox}[
    colback=white,
    colframe=black!30,
    width=0.8\textwidth,
    arc=2pt,
    boxrule=0.4pt,
    top=10pt, bottom=10pt,
  ]
    \centering\small\color{DarkGray}
    A companion technical report to the first Rust implementation of the 
    Aleph Filter, based on the VLDB 2024 paper by Dayan, Bercea, and 
    Pagh.\\[6pt]
    \texttt{arxiv.org/abs/2404.04703}
    \quad$\cdot$\quad
    \texttt{github.com/NPX2218/aleph-filter}
  \end{tcolorbox}
  \vspace{1cm}
  
\end{titlepage}

% --- TABLE OF CONTENTS ---
\pagenumbering{roman}
\tableofcontents
\newpage
\pagenumbering{arabic}


% ============================================================================
% CHAPTER 1: BACKGROUND - PROBABILISTIC FILTERS
% ============================================================================

\chapter{Background: Probabilistic Filters}

\section{What Are Filters?}

\subsection{Hash Functions and Collisions}

Let $U$ denote the universe of all possible elements, and let $S \subseteq U$ 
with $n = |S|$ denote the set of elements we wish to represent.

A \textbf{hash function} is a deterministic function $h : U \to \{0, 1, \dots, m-1\}$ 
that maps elements from a universe $U$ (e.g., strings, integers) to a fixed range of 
$m$ possible values. Determinism means that for any input $x \in U$, repeated 
evaluations of $h(x)$ always yield the same result.

Because the universe $U$ is typically much larger than the range $m$ 
(i.e., $|U| \gg m$), multiple distinct inputs can map to the same output. 
This phenomenon is called a \textbf{collision}: two elements $x \neq y$ such that 
$h(x) = h(y)$. By the pigeonhole principle, if $|U| > m$, collisions are 
unavoidable and no hash function can be injective over such a domain.

A "good" hash function distributes outputs approximately uniformly: for a 
randomly chosen $x \in U$, each value in $\{0, 1, \dots, m-1\}$ should be 
roughly equally likely. This notion was formalized by \textit{universal hashing} \cite{carter},  which guarantees that for any 
$x \neq y$, the probability of collision satisfies 
$\Pr[h(x) = h(y)] \leq 1/m$. This property is essential for the 
data structures that follow.

\subsection{The Need for Filters}

Many systems store large datasets on slow media (e.g. disk, or across a 
network). A common operation is checking whether a key $x$ exists in a set $S$ 
before performing an expensive lookup. If we could cheaply determine that 
$x \notin S$, we could skip the lookup entirely.

Storing $S$ exactly requires $\Omega(n \log |U|)$ bits for $n = |S|$ elements, 
which may be prohibitively large. \textbf{Filters} address this by trading 
exactness for space: they use far less memory by tolerating a small 
probability of error.

\begin{definitionbox}[Probabilistic Filter]
A \textbf{filter} is a compact data structure that represents a set $S$ and answers

membership queries: ``Is element $x$ in $S$?'' It can return \textbf{false negatives: never},
but may return \textbf{false positives} with some bounded probability $\epsilon$.
\end{definitionbox}


\section{Bloom Filters}

A Bloom filter, introduced by Burton Bloom in 1970 \cite{bloom}, is a 
probabilistic data structure for approximate set membership testing. Given 
a set $S$, a Bloom filter can answer queries of the form ``Is $x \in S$?'' 
with a one-sided error:

\begin{itemize}
  \item If $x \in S$, the filter always returns \textsc{yes} (no false negatives).
  \item If $x \notin S$, the filter may incorrectly return \textsc{yes} with 
        probability at most $\epsilon$ (a false positive).
\end{itemize}

This property made Bloom filters practical in real systems. For example, 
Google Chrome \cite{chrome-bloom} used a Bloom filter until 2012 to check URLs against a local 
list of known malicious sites. If the filter returned \textsc{no}, the URL 
was safe and no further lookup was needed. Only on a \textsc{yes} did the 
browser perform a more expensive server-side check.

\subsection{Structure and Operations}

A Bloom filter consists of two components
\begin{enumerate}
  \item A \textbf{bit array} $B$ of $m$ bits, where:
  \[
  B[i] \leftarrow 0 \quad \forall \, i \in \{0, 1, 2, \dots, m-1\}
\]

  \item A family of $k$ independent hash functions $h_1, h_2, \dots, h_k$, 
        each mapping elements from the universe $U$ uniformly to 
        $\{0, 1, \dots, m-1\}$.
\end{enumerate}
The Bloom filter also supports the following operations:

\begin{enumerate}
  \item \textbf{Insertion:} To insert an element $x$ into the filter, compute 
$h_1(x), h_2(x), \dots, h_k(x)$ and set each corresponding bit to 1:
\[
  B[h_i(x)] \leftarrow 1 \quad \forall \, i \in \{1, 2 \dots, k\}
\]
\item \textbf{Query:} To test whether an element $x$ is in the set, check 
whether all $k$ bits are set:
\[
  \text{query}(x) = \bigwedge_{i=1}^{k} B[h_i(x)]
\]

\end{enumerate}



If any bit $B[h_i(x)] = 0$, then $x$ was never inserted,  this is 
certain because insertion only sets bits to 1, never back to 0. If all 
$k$ bits are 1, then either $x$ was inserted, or other insertions 
happened to set those same bits, which creates a false positive.

\subsection{Worked Example}

Consider a Bloom filter with $m = 10$ bits and $k = 3$ hash functions 
$h_1, h_2, h_3$. We begin with an empty bit array:

% ---- Step 1: Empty array ----
\begin{figure}[H]
  \centering
  \begin{tikzpicture}
    \foreach \i in {0,...,9} {
      \node[bit, fill=BitZero] (b\i) at (\i*0.6, 0) {0};
      \node[font=\tiny, text=MedGray] at (\i*0.6, -0.45) {\i};
    }
    \node[lbl, anchor=east] at (-0.5, 0) {$B$:};
  \end{tikzpicture}
  \caption{Initial state: all bits set to 0.}
  \label{fig:bloom-empty}
\end{figure}

\textbf{Insert ``hello''.} Suppose our hash functions produce:
\[
  h_1(\texttt{"hello"}) = 1, \quad 
  h_2(\texttt{"hello"}) = 4, \quad 
  h_3(\texttt{"hello"}) = 9
\]
We set $B[1], B[4], B[9] \leftarrow 1$:

% ---- Step 2: Insert "hello" ----
\begin{figure}[H]
  \centering
  \begin{tikzpicture}
    \foreach \i/\val in {0/0, 1/1, 2/0, 3/0, 4/1, 5/0, 6/0, 7/0, 8/0, 9/1} {
      \ifnum\val=1
        \node[bit, fill=BitOne] (b\i) at (\i*0.6, 0) {\val};
      \else
        \node[bit, fill=BitZero] (b\i) at (\i*0.6, 0) {\val};
      \fi
      \node[font=\tiny, text=MedGray] at (\i*0.6, -0.45) {\i};
    }
    \node[lbl, anchor=east] at (-0.5, 0) {$B$:};
    
    % Arrows from element
    \node[hashbox] (elem) at (2.7, 1.5) {\texttt{"hello"}};
    \draw[arrow, AccentBlue] (elem.south) -- ++(0, -0.3) -| (b1.north);
    \draw[arrow, AccentTeal] (elem.south) -- ++(0, -0.3) -| (b4.north);
    \draw[arrow, AccentOrange] (elem.south) -- ++(0, -0.3) -| (b9.north);
  \end{tikzpicture}
  \caption{After inserting \texttt{"hello"}: bits 1, 4, and 9 are set.}
  \label{fig:bloom-insert-hello}
\end{figure}

\textbf{Insert ``world''.} Suppose:
\[
  h_1(\texttt{"world"}) = 3, \quad 
  h_2(\texttt{"world"}) = 4, \quad 
  h_3(\texttt{"world"}) = 7
\]
We set $B[3], B[4], B[7] \leftarrow 1$. Note that $B[4]$ was already 1 
from inserting \texttt{"hello"}, it remains 1. This overlap is normal 
and is the source of false positives.

% ---- Step 3: Insert "world" ----
\begin{figure}[H]
  \centering
  \begin{tikzpicture}
    \foreach \i/\val in {0/0, 1/1, 2/0, 3/1, 4/1, 5/0, 6/0, 7/1, 8/0, 9/1} {
      \ifnum\val=1
        \node[bit, fill=BitOne] (b\i) at (\i*0.6, 0) {\val};
      \else
        \node[bit, fill=BitZero] (b\i) at (\i*0.6, 0) {\val};
      \fi
      \node[font=\tiny, text=MedGray] at (\i*0.6, -0.45) {\i};
    }
    \node[lbl, anchor=east] at (-0.5, 0) {$B$:};
    
    % Arrows from element
    \node[hashbox] (elem) at (3.9, 1.5) {\texttt{"world"}};
    \draw[arrow, AccentBlue] (elem.south) -- ++(0, -0.3) -| (b3.north);
    \draw[arrow, AccentTeal] (elem.south) -- ++(0, -0.3) -| (b4.north);
    \draw[arrow, AccentOrange] (elem.south) -- ++(0, -0.3) -| (b7.north);
  \end{tikzpicture}
  \caption{After inserting \texttt{"world"}: bits 3, 4, and 7 are set. 
           Bit 4 was already set by \texttt{"hello"}.}
  \label{fig:bloom-insert-world}
\end{figure}

Now we query three elements that were \emph{never} inserted to illustrate 
the filter's behavior.

\textbf{Query ``cat'', true negative.} Suppose:
\[
  h_1(\texttt{"cat"}) = 2, \quad 
  h_2(\texttt{"cat"}) = 5, \quad 
  h_3(\texttt{"cat"}) = 7
\]
We check $B[2] = 0$. Since at least one bit is 0, the filter returns 
\textsc{no} --- \texttt{"cat"} is \textbf{definitely not} in the set. 
This answer is guaranteed correct.

% ---- Step 4: Query "cat" — true negative ----
\begin{figure}[H]
  \centering
  \begin{tikzpicture}
    \foreach \i/\val in {0/0, 1/1, 2/0, 3/1, 4/1, 5/0, 6/0, 7/1, 8/0, 9/1} {
      \ifnum\val=1
        \node[bit, fill=BitOne] (b\i) at (\i*0.6, 0) {\val};
      \else
        \node[bit, fill=BitZero] (b\i) at (\i*0.6, 0) {\val};
      \fi
      \node[font=\tiny, text=MedGray] at (\i*0.6, -0.45) {\i};
    }
    \node[lbl, anchor=east] at (-0.5, 0) {$B$:};
    
    % Query arrows
    \node[hashbox] (elem) at (2.4, 1.5) {\texttt{"cat"}};
    \draw[arrow, AccentRed] (elem.south) -- ++(0, -0.3) -| (b2.north);
    \draw[arrow, AccentRed] (elem.south) -- ++(0, -0.3) -| (b5.north);
    \draw[arrow, AccentGreen] (elem.south) -- ++(0, -0.3) -| (b7.north);
    
    % X marks on the 0-bits
    \node[font=\bfseries, text=AccentRed] at (2*0.6, 0.65) {\texttimes};
    \node[font=\bfseries, text=AccentRed] at (5*0.6, 0.65) {\texttimes};
        \node[font=\small, text=AccentGreen] at (7*0.6, 0.65) {\checkmark};

    
    % Result
    \node[lbl, text=AccentRed, anchor=west] at (6.5, 0) 
      {$B[2] = 0 \Rightarrow$ \textbf{definitely not in set}};
  \end{tikzpicture}
  \caption{Querying \texttt{"cat"}: bit 2 is 0, so the filter correctly 
           returns \textsc{no}.}
  \label{fig:bloom-query-cat}
\end{figure}

\textbf{Query ``dog'' yields a false positive.} Suppose:
\[
  h_1(\texttt{"dog"}) = 1, \quad 
  h_2(\texttt{"dog"}) = 3, \quad 
  h_3(\texttt{"dog"}) = 7
\]
We check $B[1] = 1$, $B[3] = 1$, $B[7] = 1$. All three bits are set, 
so the filter returns \textsc{yes}. But \texttt{"dog"} was never inserted, bit 1 was set by \texttt{"hello"}, bit 3 by \texttt{"world"}, and 
bit 7 by \texttt{"world"}. This is a \textbf{false positive}: the filter 
is wrong, but has no way to know it.

% ---- Step 5: Query "dog" — false positive ----
\begin{figure}[H]
  \centering
  \begin{tikzpicture}
    \foreach \i/\val in {0/0, 1/1, 2/0, 3/1, 4/1, 5/0, 6/0, 7/1, 8/0, 9/1} {
      \ifnum\val=1
        \node[bit, fill=BitOne] (b\i) at (\i*0.6, 0) {\val};
      \else
        \node[bit, fill=BitZero] (b\i) at (\i*0.6, 0) {\val};
      \fi
      \node[font=\tiny, text=MedGray] at (\i*0.6, -0.45) {\i};
    }
    \node[lbl, anchor=east] at (-0.5, 0) {$B$:};
    
    % Query arrows
    \node[hashbox] (elem) at (2.4, 1.5) {\texttt{"dog"}};
    \draw[arrow, AccentGreen] (elem.south) -- ++(0, -0.3) -| (b1.north);
    \draw[arrow, AccentGreen] (elem.south) -- ++(0, -0.3) -| (b3.north);
    \draw[arrow, AccentGreen] (elem.south) -- ++(0, -0.3) -| (b7.north);
    
    % Check marks on the 1-bits
    \node[font=\small, text=AccentGreen] at (1*0.6, 0.65) {\checkmark};
    \node[font=\small, text=AccentGreen] at (3*0.6, 0.65) {\checkmark};
    \node[font=\small, text=AccentGreen] at (7*0.6, 0.65) {\checkmark};
    
    % Result
    \node[lbl, text=AccentOrange, anchor=west, align=left] at (6.5, 0.3) 
      {All bits are 1};
    \node[lbl, text=AccentOrange, anchor=west, align=left] at (6.5, -0.2) 
      {$\Rightarrow$ \textbf{false positive!}};
  \end{tikzpicture}
  \caption{Querying \texttt{"dog"}: all checked bits happen to be 1 
           (set by other insertions), so the filter incorrectly returns 
           \textsc{yes}. This is a false positive.}
  \label{fig:bloom-query-dog}
\end{figure}

\begin{insightbox}[Why False Positives Happen]
A Bloom filter only stores bits, not the elements themselves. When multiple 
elements set overlapping bits, the filter cannot distinguish whether a 
particular combination of 1-bits was caused by the queried element or by 
a coincidence of other insertions. This is the fundamental tradeoff: we 
use far less memory than storing elements explicitly, but we accept a 
small probability $\epsilon$ of false positives.
\end{insightbox}

This example also reveals two important limitations of the Bloom filter 
that motivate the data structures in later chapters:

\begin{enumerate}
  \item \textbf{No deletion:} We cannot remove \texttt{"hello"} from the 
        filter by setting bits 1, 4, and 9 back to 0, as the bit 4 is shared 
        with \texttt{"world"}, so clearing it would create a 
        \emph{false negative}, which Bloom filters must never produce.
  \item \textbf{No expansion:} The bit array size $m$ is fixed at creation. 
        As more elements are inserted, the array fills with 1-bits, and 
        the false positive rate increases. The only remedy is to rebuild 
        the entire filter with a larger $m$, which requires access to all 
        original elements.
\end{enumerate}


\subsection{False Positive Rate}

We can quantify the probability of a false positive, using an analysis that follows the treatment in \cite{mitzenmacher}. 

\[
  \Pr[\text{bit is picked}] = \frac{1}{m}
\]

\[
  \Pr[\text{bit is not picked}] = 1 - \frac{1}{m}
\]

As inserting on element runs $k$ hash functions, the probability that a specific bit is not set by one element is:
\[
    \Pr[\text{bit not set by one element}] = \left(1 -
    \frac{1}{m}\right)^k
\]

After inserting $n$ 
elements into a Bloom filter with $m$ bits and $k$ hash functions, the 
probability that a single bit is still 0 is:
\[
  \Pr[\text{bit} = 0] = \left(1 - \frac{1}{m}\right)^{kn} = \left(\left(1 - \frac{1}{m}\right)^{m}\right)^{kn/m}
\]

When $m$ is large, we can use the approximation $(1 - 1/m)^m \approx e^{-1}$ to simplify the expression.

\[
\left(\left(1 - \frac{1}{m}\right)^{m}\right)^{kn/m} \approx (e^{-1})^{kn/m} = e^{-kn/m}
\]

A false positive occurs when all $k$ bits checked for a non-member happen 
to be 1. Since each bit is 1 with probability $1 - e^{-kn/m}$, the false 
positive rate is:

\begin{equation}
  \epsilon \approx \left(1 - e^{-kn/m}\right)^k
  \label{eq:bloom-fpr}
\end{equation}

To find the optimal number of hash functions $k$ that minimizes $\epsilon$ for a given $m/n$, we can take the derivative of $\epsilon$ with respect to $k$ and set it to zero.

Let $p = e^{-kn/m}$. Then $\epsilon$ can be rewritten as:
\[
  \epsilon = (1 - p)^k
\]
Taking the natural logarithm:
\[
    \ln \epsilon = k \ln(1 - p)
\]
Now, we differentiate with respect to $k$:
\[
    \frac{d}{dk} \ln \epsilon = \ln(1 - p
) + k \cdot \frac{-1}{1 - p} \cdot \frac{dp}{dk}
\]
Calculating $\frac{dp}{dk}$:
\[
    \frac{dp}{dk} = \frac{d}{dk} e^{-kn/m} = -\frac{n}{m} e^{-kn/m} = -\frac{n}{m} p
\]
Substituting back:
\[
    \frac{d}{dk} \ln \epsilon = \ln(1 - p) + k \cdot \frac{p}{1 - p} \cdot \frac{n}{m}
\]  Setting this derivative to zero for optimality:
\[
    \ln(1 - p) + k \cdot \frac{p}{1 - p} \cdot \frac{n}{m} = 0
\]  
Rearranging gives:
\[    k \cdot \frac{p}{1 - p} \cdot \frac{n}{m} = -\ln(1 - p) \]  


\[
    k = -\frac{m}{n}\ln(1-p) \cdot \frac{(1-p)}{p}
\]
From our definition of $p$ we get that:

\[
 k = \frac{-m\ln(p)}{n}
\]
Setting the two equation equal to each other

\[
 -\frac{m}{n}\ln(1-p) \cdot \frac{(1-p)}{p} = -\frac{m}n \cdot \ln(p)
\]
Cancel like terms and simplify:

\[
p\ln(p)  = (1-p)\ln(1-p)
\]
We get that $p = 1/2$, meaning that 

\[
 k = \frac{-m\ln(1/2)}{n} = \frac{m}{n}\ln(2)
\]
Formally we can state that, for a given ratio of bits per element $m/n$, where $m$ is the size of the bit array and $n$ is the number of 
elements to store, the optimal number of hash functions that minimizes 
$\epsilon$ is:
\begin{equation}
  k_{\text{opt}} = \frac{m}{n} \ln 2
  \label{eq:bloom-optimal-k}
\end{equation}

Substituting $k_{\text{opt}}$ back into Equation~\ref{eq:bloom-fpr} gives:
\begin{equation}
  \epsilon_{\text{opt}} \approx \left(\frac{1}{2}\right)^{k_{\text{opt}}} 
  = \left(\frac{1}{2}\right)^{\frac{m}{n} \ln 2} =  \left(\left(\frac{1}{2}\right)^{\ln 2}\right)^{\frac{m}{n}} = (0.6185)^{\frac{m}{n}}
  \label{eq:bloom-fpr-optimal}
\end{equation}

This tells us that the false positive rate decreases \emph{exponentially} 
as we allocate more bits per element. The following table and figure 
illustrate this relationship:

\begin{table}[H]
  \centering
  \caption{False positive rate vs.\ bits per element (with optimal $k$)}
  \label{tab:bloom-fpr}
  \begin{tabular}{c c c}
    \hline
    \textbf{Bits per element} ($m/n$) & \textbf{Optimal $k$} & \textbf{FPR} ($\epsilon$) \\
    \hline
    4  & 3  & 14.7\%   \\
    6  & 4  & 5.6\%   \\
    8  & 6  & 2.2\%  \\
    10 & 7  & 0.82\% \\
    12 & 8  & 0.31\% \\
    16 & 11 & 0.046\% \\
    20 & 14 & 0.0067\% \\
    \hline
  \end{tabular}
\end{table}

% ---- FPR vs bits-per-element graph ----
\begin{figure}[H]
  \centering
  \begin{tikzpicture}[scale=1]
    % Axes
    \draw[arrow] (0, 0) -- (8.5, 0) 
      node[right, font=\small] {bits per element ($m/n$)};
    \draw[arrow] (0, 0) -- (0, 5.5) 
      node[above, font=\small] {FPR ($\epsilon$)};
    
    % Y-axis labels (percentage)
    \foreach \y/\label in {0/0\%, 1/5\%, 2/10\%, 3/15\%, 4/20\%, 5/25\%} {
      \draw (-0.1, \y) -- (0.1, \y);
      \node[font=\scriptsize, anchor=east] at (-0.15, \y) {\label};
    }
    
    % X-axis labels
    \foreach \x/\label in {1/2, 2/4, 3/6, 4/8, 5/10, 6/12, 7/16, 8/20} {
      \draw (\x, -0.1) -- (\x, 0.1);
      \node[font=\scriptsize, anchor=north] at (\x, -0.15) {\label};
    }
    
    % Grid lines (light)
    \foreach \y in {1, 2, 3, 4, 5} {
      \draw[gray!20] (0, \y) -- (8.5, \y);
    }
    \foreach \x in {1, 2, 3, 4, 5, 6, 7, 8} {
      \draw[gray!20] (\x, 0) -- (\x, 5.5);
    }
    
    % Plot using corrected values:
    % y-axis: 1 unit = 5%, so y = FPR/5
    % x-axis: x=2 -> m/n=4, x=3 -> m/n=6, etc.
    %
    % m/n=4:  14.7%  -> x=2, y=2.94
    % m/n=6:  5.6%   -> x=3, y=1.12
    % m/n=8:  2.2%   -> x=4, y=0.44
    % m/n=10: 0.82%  -> x=5, y=0.164
    % m/n=12: 0.31%  -> x=6, y=0.062
    % m/n=16: 0.046% -> x=7, y=0.0092
    % m/n=20: 0.0067%-> x=8, y=0.00134
    
    % Curve
    \draw[AccentBlue, very thick, smooth] plot coordinates {
      (1, 5.5)
      (1.5, 4.2)
      (2, 2.94)
      (2.5, 1.9)
      (3, 1.12)
      (3.5, 0.72)
      (4, 0.44)
      (5, 0.164)
      (6, 0.062)
      (7, 0.0092)
      (8, 0.00134)
    };
    
    % Data points from table
    \fill[AccentBlue] (2, 2.94) circle (2pt);   % m/n=4
    \fill[AccentBlue] (3, 1.12) circle (2pt);   % m/n=6
    \fill[AccentBlue] (4, 0.44) circle (2pt);   % m/n=8
    \fill[AccentBlue] (5, 0.164) circle (2pt);  % m/n=10
    \fill[AccentBlue] (6, 0.062) circle (2pt);  % m/n=12
    \fill[AccentBlue] (7, 0.0092) circle (2pt); % m/n=16
    \fill[AccentBlue] (8, 0.00134) circle (2pt);% m/n=20
    
    % Highlighted point: 10 bits/element
    \fill[AccentRed] (5, 0.164) circle (3pt);
    \node[font=\scriptsize, text=AccentRed, anchor=south west] at (5.1, 0.3) 
      {$m/n = 10 \Rightarrow \epsilon \approx 0.82\%$};
    
    % Highlighted point: 8 bits/element  
    \fill[AccentTeal] (4, 0.44) circle (3pt);
    \node[font=\scriptsize, text=AccentTeal, anchor=south west] at (4.1, 0.6) 
      {$m/n = 8 \Rightarrow \epsilon \approx 2.2\%$};
    
    % Annotation: diminishing returns region
    \draw[brace] (5, -1) -- (8, -1) 
      node[midway, below=0pt, font=\scriptsize, text=DarkGray, align=center] 
      {diminishing returns:\\FPR already very low};
    
    % Annotation: steep decline region
    \draw[brace] (1, -1) -- (4, -1) 
      node[midway, below=0pt, font=\scriptsize, text=DarkGray, align=center] 
      {adding bits here\\has the biggest impact};
  \end{tikzpicture}
  \caption{False positive rate as a function of bits per element with 
           optimal $k$. The rate decreases exponentially: doubling the 
           bits per element roughly squares the improvement. This means that as the total ratio of bits per element increases, the FPR goes down.}
  \label{fig:bloom-fpr-curve}
\end{figure}

Equation~\ref{eq:bloom-fpr-optimal} and Figure~\ref{fig:bloom-fpr-curve} 
reveal an important practical insight: a Bloom filter with just 10 bits per 
element achieves a false positive rate under 1\%, regardless of how large 
the universe $U$ is. This extreme space efficiency, requiring only 
$\approx 1.25$ bytes per element for sub-1\% error, is what makes Bloom 
filters so widely deployed.

\subsection{Space Efficiency}

From an information-theoretic perspective, any data structure that answers 
membership queries with a false positive rate of $\epsilon$ must use at 
least $\log_2(1/\epsilon)$ bits per element, this is the minimum amount 
of information needed to distinguish ``in the set'' from ``not in the set'' 
with error probability $\epsilon$.

From Equation~\ref{eq:bloom-fpr-optimal}, we can solve for the number of 
bits per element required by a Bloom filter. Taking the logarithm of both 
sides gives:
\[
  \frac{m}{n} = \frac{\log_2(1/\epsilon)}{\ln 2} = 1.44 \log_2(1/\epsilon)
\]

A space-optimized Bloom filter therefore uses $1.44 \log_2(1/\epsilon)$ bits 
per element \cite{bloom, cuckoo}, which is a 44\% overhead over the 
information-theoretic lower bound of $\log_2(1/\epsilon)$ bits. In practice, 
this overhead is a modest price for the simplicity and speed of the data 
structure.

\begin{comment}
% ---- EXAMPLE: Bloom Filter bit array diagram ----
\begin{figure}[H]
  \centering
  \begin{tikzpicture}
    % Draw bit array
    \foreach \i/\val in {0/0, 1/1, 2/0, 3/1, 4/1, 5/0, 6/0, 7/1, 8/0, 9/1} {
      \ifnum\val=1
        \node[bit, fill=BitOne] (b\i) at (\i*0.6, 0) {\val};
      \else
        \node[bit, fill=BitZero] (b\i) at (\i*0.6, 0) {\val};
      \fi
      \node[font=\tiny, text=MedGray] at (\i*0.6, -0.45) {\i};
    }
    
    % Label
    \node[lbl, anchor=east] at (-0.5, 0) {Bit array:};
    
    % Show hash arrows from an element
    \node[hashbox] (elem) at (2.7, 1.5) {$h(x)$};
    \draw[arrow, AccentBlue] (elem.south) -- ++(0, -0.3) -| (b1.north);
    \draw[arrow, AccentTeal] (elem.south) -- ++(0, -0.3) -| (b4.north);
    \draw[arrow, AccentOrange] (elem.south) -- ++(0, -0.3) -| (b9.north);
    
    % Annotation
    \node[lbl, anchor=west] at (6.5, 1.5) {Element $x$ hashed};
    \node[lbl, anchor=west] at (6.5, 1.0) {to positions 1, 4, 9};
  \end{tikzpicture}
  \caption{A Bloom filter with 10 bits and 3 hash functions.}
  \label{fig:bloom-basic}
\end{figure}
\end{comment}


\section{Cuckoo Filters}

Cuckoo hashing is a collision resolution scheme that guarantees worst-case 
constant-time lookups. While several variants exist, we present the standard 
formulation using two hash tables.

A cuckoo hash table consists of two tables $T_1$ and $T_2$, each of size $r$, 
along with two independent hash functions $h_1, h_2 : U \to \{0, 1, \ldots, r-1\}$. 
Each element $x \in S$ is stored in exactly one of two possible locations: 
either $T_1[h_1(x)]$ or $T_2[h_2(x)]$.


\subsection{Structure and Operations}

The data structure is initalized to 
\[
  T_1[j] \leftarrow \varnothing \quad \text{and} \quad T_2[j] \leftarrow \varnothing 
  \quad \forall \, j \in \{0, 1, \dots, rr-1\}
\]
The data structure supports three operations:
\begin{enumerate}
  \item \textbf{Search:} To determine whether $x \in S$, we examine 
        $T_1[h_1(x)]$ and $T_2[h_2(x)]$. Since each lookup is a direct 
        array access, search runs in $O(1)$ worst-case time.
        
  \item \textbf{Delete:} To remove an element $x$, we check both candidate 
        positions $T_1[h_1(x)]$ and $T_2[h_2(x)]$. If $x$ occupies either 
        location, we remove it. This operation also runs in $O(1)$ worst-case time.
        
  \item \textbf{Insert:} To insert an element $x$, we attempt to place it 
        in $T_1[h_1(x)]$. If this position is occupied by some element $y$, 
        we evict $y$ and attempt to place $y$ in its alternate location. 
        This process continues until either an empty slot is found or a 
        cycle is detected, necessitating a rehash with new hash functions.
\end{enumerate}



\subsection{Worked Example}

We begin with two empty tables of size 4:

% ---- Step 1: Empty tables ----
\begin{figure}[H]
  \centering
  \begin{tikzpicture}
    % Table 1
    \node[lbl, font=\bfseries] at (1.95, 1) {$T_1$ (via $h_1$)};
    \foreach \i in {0,...,3} {
      \node[emptyslot, minimum width=1.3cm] (a\i) at (\i*1.4, 0.3) {};
      \node[font=\tiny, text=MedGray] at (\i*1.4, -0.25) {\i};
    }
    
    % Table 2
    \node[lbl, font=\bfseries] at (1.95, -1.2) {$T_2$ (via $h_2$)};
    \foreach \i in {0,...,3} {
      \node[emptyslot, minimum width=1.3cm] (b\i) at (\i*1.4, -1.9) {};
      \node[font=\tiny, text=MedGray] at (\i*1.4, -2.45) {\i};
    }
  \end{tikzpicture}
  \caption{Two empty hash tables for cuckoo hashing.}
  \label{fig:cuckoo-empty}
\end{figure}

\textbf{Insert ``A'':} Suppose $h_1(\texttt{A}) = 1$ and $h_2(\texttt{A}) = 2$.
Slot 1 in $T_1$ is empty, so we place \texttt{A} there:

% ---- Step 2: Insert A ----
\begin{figure}[H]
  \centering
  \begin{tikzpicture}
    % Table 1
    \node[lbl, font=\bfseries] at (1.95, 1) {$T_1$};
    \foreach \i in {0,...,3} {
      \node[emptyslot, minimum width=1.3cm] (a\i) at (\i*1.4, 0.3) {};
      \node[font=\tiny, text=MedGray] at (\i*1.4, -0.25) {\i};
    }
    \node[fullslot, minimum width=1.3cm] at (1*1.4, 0.3) {\fp{A}};
    
    % Table 2
    \node[lbl, font=\bfseries] at (1.95, -1.2) {$T_2$};
    \foreach \i in {0,...,3} {
      \node[emptyslot, minimum width=1.3cm] (b\i) at (\i*1.4, -1.9) {};
      \node[font=\tiny, text=MedGray] at (\i*1.4, -2.45) {\i};
    }
    
    % Arrow showing where A went
    \node[hashbox] (key) at (-1.8, 0.3) {\texttt{A}};
    \draw[arrow, AccentBlue] (key.east) -- (a1.west)
      node[midway, above=10pt, font=\scriptsize, text=AccentBlue] {$h_1 = 1$};
  \end{tikzpicture}
  \caption{Insert \texttt{A}: placed in $T_1[1]$.}
  \label{fig:cuckoo-insert-a}
\end{figure}

\textbf{Insert ``B'':} Suppose $h_1(\texttt{B}) = 3$ and $h_2(\texttt{B}) = 0$.
Slot 3 in $T_1$ is empty, so we place \texttt{B} there:

% ---- Step 3: Insert B ----
\begin{figure}[H]
  \centering
  \begin{tikzpicture}
    % Table 1
    \node[lbl, font=\bfseries] at (1.95, 1) {$T_1$};
    \foreach \i in {0,...,3} {
      \node[emptyslot, minimum width=1.3cm] (a\i) at (\i*1.4, 0.3) {};
      \node[font=\tiny, text=MedGray] at (\i*1.4, -0.25) {\i};
    }
    \node[fullslot, minimum width=1.3cm] at (1*1.4, 0.3) {\fp{A}};
    \node[fullslot, minimum width=1.3cm] at (3*1.4, 0.3) {\fp{B}};
    
    % Table 2
    \node[lbl, font=\bfseries] at (1.95, -1.2) {$T_2$};
    \foreach \i in {0,...,3} {
      \node[emptyslot, minimum width=1.3cm] (b\i) at (\i*1.4, -1.9) {};
      \node[font=\tiny, text=MedGray] at (\i*1.4, -2.45) {\i};
    }
    
    % Arrow
    \node[hashbox] (key) at (-1.8, 0.3) {\texttt{B}};
    \draw[arrow, AccentBlue] (key.east) -- (a3.west)
      node[midway, above=10pt, font=\scriptsize, text=AccentBlue, xshift=-10px] {$h_1 = 3$};
  \end{tikzpicture}
  \caption{Insert \texttt{B}: placed in $T_1[3]$.}
  \label{fig:cuckoo-insert-b}
\end{figure}

\textbf{Insert ``C'':} Suppose $h_1(\texttt{C}) = 1$ and 
$h_2(\texttt{C}) = 3$. Slot 1 in $T_1$ is occupied by \texttt{A}. So 
\texttt{C} \emph{evicts} \texttt{A} from $T_1[1]$, and \texttt{A} moves 
to its alternative location $h_2(\texttt{A}) = 2$ in $T_2$:

% ---- Step 4: Insert C with eviction ----
\begin{figure}[H]
  \centering
  \begin{tikzpicture}
    % Table 1
    \node[lbl, font=\bfseries] at (1.95, 1) {$T_1$};
    \foreach \i in {0,...,3} {
      \node[emptyslot, minimum width=1.3cm] (a\i) at (\i*1.4, 0.3) {};
      \node[font=\tiny, text=MedGray] at (\i*1.4, -0.25) {\i};
    }
    \node[fullslot, minimum width=1.3cm] (ac) at (1*1.4, 0.3) {\fp{C}};
    \node[fullslot, minimum width=1.3cm] at (3*1.4, 0.3) {\fp{B}};
    
    % Table 2
    \node[lbl, font=\bfseries] at (1.95, -1.2) {$T_2$};
    \foreach \i in {0,...,3} {
      \node[emptyslot, minimum width=1.3cm] (b\i) at (\i*1.4, -1.9) {};
      \node[font=\tiny, text=MedGray] at (\i*1.4, -2.45) {\i};
    }
    \node[fullslot, minimum width=1.3cm] (ba) at (2*1.4, -1.9) {\fp{A}};
    
    % Eviction arrow
    \draw[arrow, AccentRed, thick] (ac.south) to[out=-100, in=90] (ba.north)
      node[midway, yshift=-25px, right=4pt, font=\scriptsize, text=AccentRed, align=left] 
      {\texttt{A} evicted\\to $T_2[2]$};
    
    % Insert arrow
    \node[hashbox] (key) at (-1.8, 0.3) {\texttt{C}};
    \draw[arrow, AccentBlue] (key.east) -- (ac.west)
      node[midway, above=10pt, font=\scriptsize, text=AccentBlue] {$h_1 = 1$};
  \end{tikzpicture}
  \caption{Insert \texttt{C}: evicts \texttt{A} from $T_1[1]$. 
           \texttt{A} moves to its alternative slot $T_2[2]$.}
  \label{fig:cuckoo-eviction}
\end{figure}

\textbf{Query ``A'': } To look up a key, we check both possible locations:
$T_1[h_1(\texttt{A})] = T_1[1]$ contains \texttt{C} (not a match), but 
$T_2[h_2(\texttt{A})] = T_2[2]$ contains \texttt{A}, meaning it was found. This 
always takes at most 2 lookups, giving $\bigO{1}$ query time.

\textbf{Delete ``A'':} To delete, we simply clear $T_2[2]$. Unlike Bloom 
filters, deletion is straightforward because each key occupies exactly 
one slot, clearing it cannot affect other keys.

\begin{insightbox}[Bloom Filters vs.\ Cuckoo Hashing]
In a Bloom filter, multiple keys share the same bits, making deletion 
impossible without risking false negatives. In cuckoo hashing, each key 
has its own dedicated slot, so deletion is safe. However, the table has 
a fixed capacity, meaning that once both tables are sufficiently full, insertions 
fail and the entire structure must be rebuilt. This inability to 
\emph{expand} efficiently is the limitation that quotient filters and 
ultimately the Aleph Filter address.
\end{insightbox}


% ---- EXAMPLE: Cuckoo filter with two possible slots ----
\begin{comment}
\begin{figure}[H]
  \centering
  \begin{tikzpicture}
    % Bucket array
    \foreach \i in {0,...,7} {
      \node[emptyslot] (s\i) at (\i*1.3, 0) {};
      \node[font=\tiny, text=MedGray] at (\i*1.3, -0.55) {\i};
    }
    
    % Fill some slots with fingerprints
    \node[fullslot] at (1*1.3, 0) {\fp{0xA3}};
    \node[fullslot] at (3*1.3, 0) {\fp{0x7F}};
    \node[fullslot] at (5*1.3, 0) {\fp{0x12}};
    
    % Show element hashing to two buckets
    \node[above=1.2cm of s2, font=\small] (key) {key $k$};
    \draw[arrow, AccentBlue] (key.south) -- (s2.north) 
      node[midway, left, font=\scriptsize, text=AccentBlue] {$h_1(k)$};
    \draw[arrow, AccentOrange] (key.south) to[out=-20, in=90] (s6.north)
      node[midway, right, font=\scriptsize, text=AccentOrange] {$h_2(k)$};
    
    \node[lbl] at (4.5, -1.2) {Two candidate buckets per key};
  \end{tikzpicture}
  \caption{Cuckoo filter: each key maps to two possible buckets via two hash functions.}
  \label{fig:cuckoo-basic}
\end{figure}

\end{comment}
% ============================================================================
% FROM CUCKOO HASHING TO CUCKOO FILTERS — formal version with diagrams
% Paste this after your cuckoo hashing worked example
% ============================================================================
% ============================================================================
% FROM CUCKOO HASHING TO CUCKOO FILTERS — clean version, no repeats
% ============================================================================

\subsection{From Cuckoo Hashing to Cuckoo Filters}

The cuckoo hashing scheme presented above stores complete keys in each 
table entry. A \textbf{cuckoo filter} \cite{cuckoo} modifies this design 
by storing only an $f$-bit \textbf{fingerprint} $\phi(x)$ for each 
element $x$, where $\phi : U \to \{0,1\}^f$ is a hash function mapping 
elements to compact bit strings. This transformation converts an exact 
hash table into a probabilistic filter: it uses significantly less space, 
but introduces false positives when two distinct elements $x \neq y$ 
happen to satisfy $\phi(x) = \phi(y)$.

\begin{definitionbox}[Fingerprint]
A \textbf{fingerprint} $\phi(x)$ is an $f$-bit hash of element $x$. 
The parameter $f$ controls the tradeoff between space and accuracy: 
larger $f$ means more possible fingerprint values ($2^f$), reducing the 
chance of collisions but increasing memory usage.
\end{definitionbox}

\subsubsection{Single-Array Structure and Bucket Positions}

While our earlier presentation of cuckoo hashing used two separate 
tables $T_1$ and $T_2$, the cuckoo filter consolidates them into a 
single array $B$ of $m$ buckets using \textbf{partial-key cuckoo hashing} 
\cite{cuckoo}. To understand how the two candidate buckets are computed, 
consider inserting the element \texttt{"hello"}:

\begin{enumerate}
  \item \textbf{Hash the full key:} Compute a hash of the original 
        element, producing a large integer:
        \[
          H(\texttt{"hello"}) = \texttt{0xA3F7B21E9C4D8256} \quad \text{(64 bits)}
        \]
  
  \item \textbf{Split the hash:} Divide the output into two 
        non-overlapping parts:
        \begin{itemize}
          \item The \textbf{long part} determines the first bucket index.
          \item The \textbf{short part} becomes the fingerprint 
                $\phi(x)$, the only piece stored in the table.
        \end{itemize}
        
% ---- Diagram: Hash splitting ----
\begin{figure}[H]
  \centering
  \begin{tikzpicture}
    \node[font=\small] at (-2.2, 0) {$H(x) =$};
    
    % Long part
    \foreach \i/\val in {0/A, 1/3, 2/F, 3/7, 4/B, 5/2, 6/1, 7/E, 
                          8/9, 9/C, 10/4, 11/D, 12/8, 13/2} {
      \node[bit, fill=AccentTeal!15, minimum width=0.42cm] at (\i*0.42, 0) 
        {\tiny\texttt{\val}};
    }
    
    % Short part (fingerprint)
    \foreach \i/\val in {14/5, 15/6} {
      \node[bit, fill=AccentOrange!25, minimum width=0.42cm] at (\i*0.42, 0) 
        {\tiny\texttt{\val}};
    }
    
    % Braces
    \draw[brace] (0*0.42-0.18, 0.45) -- (13*0.42+0.18, 0.45)
      node[midway, above=8pt, font=\scriptsize, text=AccentTeal] 
      {used to compute $h_1$ (then discarded)};
    \draw[brace] (14*0.42-0.18, 0.45) -- (15*0.42+0.18, 0.45)
      node[midway, above=20pt, font=\scriptsize, text=AccentOrange] 
      {fingerprint $\phi(x)$ (stored)};
  \end{tikzpicture}
  \caption{A single hash of the key is split: the long part computes 
           the bucket index, the short part becomes the stored fingerprint.
           The two parts must not overlap to avoid correlation.}
  \label{fig:hash-split}
\end{figure}

  \item \textbf{Compute the first bucket $h_1$:} Apply the modulo 
        operation to the long part. The \textbf{modulo} operation 
        $a \mod m$ returns the remainder after dividing $a$ by $m$, 
        guaranteeing the result falls in $\{0, 1, \dots, m-1\}$ 
        regardless of how large the hash value is:
        \[
          h_1(x) = \texttt{0xA3F7B21E9C4D82} \mod m
        \]
        For example, with $m = 8$ buckets:
        \[
          \texttt{0xA3F7B21E9C4D82} \mod 8 = 2
        \]
        So $h_1(x) = 2$.

  \item \textbf{Compute the second bucket $h_2$:} Hash the fingerprint 
        with a separate hash function, then XOR the result with $h_1$. 
        Suppose $G(\texttt{0x56}) \mod m = 4$:
        \[
        h_2(x) = h_1(x) \oplus G(\phi(x)) = 2 \oplus 4 = 6
        \]
\end{enumerate}

\noindent In general, for any element $x$:
\begin{align}
  h_1(x) &= H(x) \mod m \label{eq:cuckoo-h1} \\
h_2(x) &= h_1(x) \oplus G(\phi(x)) \label{eq:cuckoo-h2}
\end{align}

After insertion, only the fingerprint (\texttt{0x56}) is stored in one 
of the two candidate buckets. The original key \texttt{"hello"} and the 
long part of the hash are discarded permanently.

% ---- Diagram: Single array with two candidate buckets ----
\begin{figure}[H]
  \centering
  \begin{tikzpicture}
    % Array of buckets
    \foreach \i in {0,...,7} {
      \node[emptyslot, minimum width=1.1cm] (s\i) at (\i*1.2, 0) {};
      \node[font=\tiny, text=MedGray] at (\i*1.2, -0.55) {\i};
    }
    
    % Some occupied slots
    \node[fullslot, minimum width=1.1cm] at (0*1.2, 0) {\fp{0xA3}};
    \node[fullslot, minimum width=1.1cm] at (3*1.2, 0) {\fp{0x7F}};
    \node[fullslot, minimum width=1.1cm] at (5*1.2, 0) {\fp{0x12}};
    
    % Element being inserted
    \node[hashbox] (key) at (3.6, 2.2) {\texttt{"hello"}};
    
    % Fingerprint
    \node[font=\scriptsize, text=DarkGray, xshift=50px, yshift=2.5px] at (3.6, 1.5) 
      {$\phi(x) = \texttt{0x56}$};
    
    % Two candidate buckets
    \draw[arrow, AccentBlue, thick] (key.south) -- ++(0, -0.5) -| (s2.north)
      node[pos=0.75, left, font=\scriptsize, text=AccentBlue] 
      {$h_1 = 2$};
    \draw[arrow, AccentOrange, thick] (key.south) -- ++(0, -0.5) -| (s6.north)
      node[pos=0.75, right, font=\scriptsize, text=AccentOrange] 
      {$h_2 = 6$};
    
    % Label
    \node[lbl, font=\small] at (4.2, -1.2) 
      {Single array $B$: fingerprint \texttt{0x56} can go in bucket 2 or 6};
  \end{tikzpicture}
  \caption{The cuckoo filter stores fingerprints in a single array. 
           Element \texttt{"hello"} maps to buckets 2 and 6; only the 
           fingerprint \texttt{0x56} is stored in one of them.}
  \label{fig:cuckoo-filter-array}
\end{figure}

\begin{insightbox}[Why Re-Hash the Fingerprint?]
In step 4, the fingerprint is hashed before XOR-ing with $h_1$. Without 
this, $h_2 = h_1 \oplus \phi(x)$ directly would constrain the alternate 
bucket to within $2^f$ positions of $h_1$ (e.g., within 256 positions for 
$f = 8$). Hashing the fingerprint first spreads the alternate buckets 
uniformly across the entire table \cite{cuckoo}.
\end{insightbox}

\subsubsection{XOR Symmetry and Eviction}

The choice of XOR in Equation~\ref{eq:cuckoo-h2} is not arbitrary. XOR 
has a unique algebraic property: \textbf{it is its own inverse}. For any 
values $a$ and $b$:
\begin{equation}
  a \oplus b = c \quad \Longrightarrow \quad c \oplus b = a
  \label{eq:xor-inverse}
\end{equation}

At the bit level, XOR flips specific bits. Applying it again flips them 
back, restoring the original value:
\[
  \underset{2}{\texttt{010}} \;\oplus\; \underset{4}{\texttt{100}} 
  \;=\; \underset{6}{\texttt{110}}
  \qquad \text{and} \qquad
  \underset{6}{\texttt{110}} \;\oplus\; \underset{4}{\texttt{100}} 
  \;=\; \underset{2}{\texttt{010}}
\]

This means that given \emph{either} bucket index and the hashed 
fingerprint, the \emph{other} bucket is always recoverable:
\[
  \underbrace{h_1(x)}_{2} \;\oplus\; 
  \underbrace{\text{hash}(\phi(x))}_{4} \;=\; 
  \underbrace{h_2(x)}_{6}
  \qquad \text{and} \qquad
  \underbrace{h_2(x)}_{6} \;\oplus\; 
  \underbrace{\text{hash}(\phi(x))}_{4} \;=\; 
  \underbrace{h_1(x)}_{2}
\]

During eviction, the filter holds a fingerprint in some bucket $i$ but 
does not know whether $i$ is $h_1$ or $h_2$, it does not matter. The 
same formula always yields the alternate bucket:
\begin{equation}
  j = i \oplus \text{hash}(\phi(x))
  \label{eq:cuckoo-alternate}
\end{equation}

This is essential because the original key $x$ has been discarded. 
Without XOR symmetry, the filter would need to store the full key to 
recompute $h_1(x) = \text{hash}(x) \mod m$, defeating the purpose of 
using fingerprints.

% ---- Diagram: XOR symmetry ----
\begin{figure}[H]
  \centering
  \begin{tikzpicture}
    % Two buckets
    \node[fullslot, minimum width=2cm] (b1) at (0, 0) {\fp{0x56}};
    \node[lbl, font=\small] at (0, -0.6) {bucket $i = 2$};
    
    \node[emptyslot, minimum width=2cm] (b2) at (6, 0) {};
    \node[lbl, font=\small] at (6, -0.6) {bucket $j = 6$};
    
    % Forward arrow
    \draw[arrow, AccentBlue, thick] 
      ([yshift=3pt]b1.east) -- ([yshift=3pt]b2.west)
      node[midway, above=6pt, font=\small, text=AccentBlue] 
      {$j = 2 \oplus \text{hash}(\texttt{0x56}) = 6$};
    
    % Reverse arrow
    \draw[arrow, AccentOrange, thick] 
      ([yshift=-3pt]b2.west) -- ([yshift=-3pt]b1.east)
      node[midway, below=6pt, font=\small, text=AccentOrange] 
      {$i = 6 \oplus \text{hash}(\texttt{0x56}) = 2$};
    
    % Annotation
    \node[lbl, font=\small, align=center] at (3, -1.5) 
      {Same formula, either direction, only needs 
       fingerprint + current bucket};
  \end{tikzpicture}
  \caption{XOR symmetry enables eviction without the original key: 
           given any bucket index and the fingerprint, the alternate 
           bucket is always computable.}
  \label{fig:cuckoo-xor}
\end{figure}

\subsubsection{Bucket Structure}

Each bucket holds up to $b$ fingerprints (typically $b = 4$). An 
insertion requires eviction only when all $2b$ slots across both 
candidate buckets are occupied. With $b = 4$, the table achieves a 
\textbf{load factor} of $\alpha \approx 0.955$ before insertions begin 
to fail \cite{cuckoo}. The load factor $\alpha$ is the fraction of 
slots currently occupied, $\alpha = 0.955$ means 95.5\% of all slots 
in the table are filled.

% ---- Diagram: Bucket internals ----
\begin{figure}[H]
  \centering
  \begin{tikzpicture}
    % Bucket label
    \node[lbl, font=\bfseries] at (-1.5, 0) {Bucket $i$:};
    
    % Four slots within one bucket
    \node[fullslot, minimum width=1.5cm] (e0) at (0.3, 0) {\fp{0xA3}};
    \node[fullslot, minimum width=1.5cm] (e1) at (1.9, 0) {\fp{0x7F}};
    \node[fullslot, minimum width=1.5cm] (e2) at (3.5, 0) {\fp{0x12}};
    \node[emptyslot, minimum width=1.5cm] (e3) at (5.1, 0) {};
    
    % Brace
    \draw[brace] (e0.north west) -- (e3.north east)
      node[midway, above=10pt, font=\small, text=DarkGray] 
      {$b = 4$ fingerprint slots per bucket};
    
    % Slot labels
    \node[font=\tiny, text=MedGray] at (0.3, -0.55) {slot 0};
    \node[font=\tiny, text=MedGray] at (1.9, -0.55) {slot 1};
    \node[font=\tiny, text=MedGray] at (3.5, -0.55) {slot 2};
    \node[font=\tiny, text=MedGray] at (5.1, -0.55) {slot 3};
    
    % Arrow to empty slot
    \draw[arrow, AccentGreen, thick] (5.1, -1.9) -- ([yshift=-15pt]e3.south)
      node[pos=0, below, font=\scriptsize, text=AccentGreen] 
      {insert here (no eviction needed)};
  \end{tikzpicture}
  \caption{A single bucket with $b = 4$ slots. Three fingerprints are 
           stored; one slot remains empty. An insertion into this bucket 
           takes the empty slot without triggering eviction.}
  \label{fig:cuckoo-bucket}
\end{figure}

\subsubsection{Operations}

The cuckoo filter supports three operations, each in $\bigO{1}$ time:

\begin{enumerate}
  \item \textbf{Insertion:} To insert element $x$, compute $\phi(x)$, $h_1(x)$, 
and $h_2(x)$. If either bucket has an empty slot, place $\phi(x)$ there. 
If both buckets are full (all $2b$ slots occupied), evict a random 
fingerprint from one bucket and relocate it to its alternate bucket via 
Equation~\ref{eq:cuckoo-alternate}. This process repeats until an empty 
slot is found or a maximum number of evictions is reached.
\item \textbf{Query:} To test whether $x \in S$, compute $\phi(x)$, $h_1(x)$, 
and $h_2(x)$, then check both candidate buckets:
\[
  \text{query}(x) = \big(\phi(x) \in B[h_1(x)]\big) \;\lor\; 
                     \big(\phi(x) \in B[h_2(x)]\big)
\]
If no matching fingerprint is found, $x$ is \textbf{definitely not} in 
the set. If a match is found, $x$ is \textbf{probably} in the set, as the match may be a false positive caused by a different element with the 
same fingerprint.
\item \textbf{Deletion:} To remove $x$, locate $\phi(x)$ in either 
$B[h_1(x)]$ or $B[h_2(x)]$ and remove one copy. This is safe because 
each fingerprint occupies its own dedicated slot, removing it cannot 
affect other elements. This is a key advantage over Bloom filters, where 
bits are shared and deletion would risk creating false negatives.
\end{enumerate}

\subsubsection{False Positive Rate}

A false positive occurs when a non-member $x$ has a fingerprint matching 
one already stored in a candidate bucket. A query checks two buckets of 
$b$ entries each, at most $2b$ fingerprints. Each $f$-bit fingerprint 
matches with probability $1/2^f$, so:
\begin{equation}
  \epsilon \approx \frac{2b}{2^f}
  \label{eq:cuckoo-fpr}
\end{equation}

Solving for the minimum fingerprint size $f$:
\[
  2^f \geq \frac{2b}{\epsilon} \quad \Longrightarrow \quad 
  f \geq \log_2(1/\epsilon) + \log_2(2b)
\]

With $b = 4$, $\log_2(2 \cdot 4) = 3$. Applying the semi-sorting 
optimization from \cite{cuckoo}, which saves one bit per fingerprint by 
exploiting the irrelevance of fingerprint ordering within a bucket:
\begin{equation}
  f = \log_2(1/\epsilon) + 2
  \label{eq:cuckoo-fp-size}
\end{equation}

\subsubsection{Space Efficiency}

Since not every slot is occupied, the amortized cost per stored item 
is the fingerprint size divided by the load factor $\alpha$:
\begin{equation}
  C_{\text{cuckoo}} = \frac{f}{\alpha} 
  = \frac{\log_2(1/\epsilon) + 2}{\alpha} 
  \quad \text{bits per item}
  \label{eq:cuckoo-space}
\end{equation}

\begin{examplebox}[Space Comparison at $\epsilon = 1\%$]
\textbf{Cuckoo filter} ($b = 4$, $\alpha = 0.955$):
\[
  C_{\text{cuckoo}} = \frac{\log_2(100) + 2}{0.955} 
  = \frac{6.64 + 2}{0.955} \approx 9.1 \text{ bits/item}
\]
\textbf{Bloom filter} (Equation~\ref{eq:bloom-fpr-optimal}):
\[
  C_{\text{Bloom}} = 1.44 \cdot \log_2(100) 
  = 1.44 \times 6.64 \approx 9.6 \text{ bits/item}
\]
At $\epsilon = 1\%$, the cuckoo filter uses 5\% less space while also 
supporting deletion.
\end{examplebox}

\begin{table}[H]
  \centering
  \caption{Space comparison: Cuckoo filter ($b = 4$, $\alpha = 0.955$) 
           vs.\ Bloom filter at various false positive rates.}
  \label{tab:cuckoo-vs-bloom}
  \begin{tabular}{c c c c}
    \hline
    \textbf{FPR} ($\epsilon$) & 
    \textbf{Cuckoo} (bits/item) & 
    \textbf{Bloom} (bits/item) & 
    \textbf{Winner} \\
    \hline
    10\%     & 5.57  & 4.78  & Bloom \\
    3\%      & 7.39  & 7.29  & $\approx$ tie \\
    1\%      & 9.05  & 9.57  & Cuckoo \\
    0.1\%    & 12.53 & 14.35 & Cuckoo \\
    0.01\%   & 16.01 & 19.13 & Cuckoo \\
    \hline
  \end{tabular}
\end{table}

The crossover occurs at $\epsilon \approx 3\%$: below this threshold, 
the cuckoo filter's fixed $+2$ overhead 
(Equation~\ref{eq:cuckoo-fp-size}) becomes a diminishing fraction of the 
growing $\log_2(1/\epsilon)$ term, while the Bloom filter's $1.44\times$ 
multiplicative factor scales proportionally.

\section{Summary of Filter Limitations}

The cuckoo filter solves the Bloom filter's deletion problem but 
inherits a critical limitation: \textbf{it cannot expand}. The array 
size $m$ is fixed at construction. When the load factor approaches 
$\alpha$, insertions fail. Expanding to a larger table of size $m' > m$ 
would require recomputing:
\[
  h_1(x) = \text{hash}(x) \mod m'
\]
for every stored element. But the original keys have been discarded,
only fingerprints remain, and $\phi(x)$ does not contain enough 
information to reconstruct $\text{hash}(x)$.

\begin{warningbox}[The Expansion Problem]
Both Bloom filters and cuckoo filters have fixed capacity. The Bloom 
filter cannot delete; the cuckoo filter cannot expand. A filter that 
supports \emph{both} deletion and efficient expansion requires a 
fundamentally different approach to storing fingerprints, one where 
expansion does not require access to the original keys. This is the 
problem that quotient filters address, and that the Aleph Filter 
ultimately solves with $\bigO{1}$ operations.
\end{warningbox}


\begin{table}[H]
  \centering
  \caption{Capabilities and limitations of the filters discussed so far. 
           The gaps in this table motivate the quotient filter family 
           and the Aleph Filter.}
  \label{tab:filter-limitations}
  \begin{tabular}{l c c c}
    \hline
    \textbf{Property} & \textbf{Bloom} & \textbf{Cuckoo} & 
    \textbf{Needed} \\
    \hline
    Insert         & $\bigO{k}$    & $\bigO{1}$*   & $\bigO{1}$ \\
    Query          & $\bigO{k}$    & $\bigO{1}$    & $\bigO{1}$ \\
    Delete         & \texttimes     & \checkmark     & \checkmark \\
    Expand         & \texttimes     & \texttimes     & \checkmark \\
    Stable FPR     & \texttimes     & \checkmark     & \checkmark \\
    \hline
  \end{tabular}
  
  \smallskip
  {\small * = amortized. $k$ = number of hash functions.}
\end{table}


\begin{insightbox}[The Core Limitation]
Bloom filters can't delete. Cuckoo filters can't expand efficiently.
Quotient filters \emph{can} expand, but their query time degrades to
$\bigO{\log N}$ as void entries accumulate. The Aleph Filter solves all three.
\end{insightbox}


% ============================================================================
% CHAPTER 2: QUOTIENT FILTERS & EXPANSION
% ============================================================================
% ============================================================================
% CHAPTER 2: QUOTIENT FILTERS AND EXPANSION
% ============================================================================

\chapter{Quotient Filters and Expansion}

The previous chapter established that Bloom filters cannot delete and 
cuckoo filters cannot expand. Both limitations stem from the same root 
cause: the stored representation (bits or fingerprints) does not retain 
enough information to reconstruct the element's position in a resized 
table.

Quotient filters, introduced by Bender et al.\ \cite{quotient}, solve 
the expansion problem with an elegant insight: instead of computing the 
bucket index and fingerprint from \emph{separate} parts of the hash 
(as cuckoo filters do), they derive both from a \emph{single} hash by 
splitting it at a chosen bit position. This means the bucket index and 
fingerprint together always reconstruct the full hash, and expansion 
simply moves the split point.

\section{The Quotienting Idea}

\begin{definitionbox}[Quotienting]
Given a $p$-bit hash $h(x)$, we split it into two parts at bit position 
$q$:
\begin{itemize}
  \item \textbf{Quotient}: the top $q$ bits, denoted $h_q(x)$ 
        $\rightarrow$ determines the \emph{slot address}
  \item \textbf{Remainder}: the bottom $r = p - q$ bits, denoted $h_r(x)$ 
        $\rightarrow$ stored as the \emph{fingerprint}
\end{itemize}
The table has $2^q$ slots. Together, the quotient and remainder 
reconstruct the full hash: $h(x) = h_q(x) \| h_r(x)$, where $\|$ 
denotes concatenation.
\end{definitionbox}

This clean reconstruction property is fundamentally different from the cuckoo filter's approach. 
Recall that in a cuckoo filter, the bucket index comes from 
$H(x) \mod m$ and the fingerprint comes from a separate 
portion of the hash. The modulo operation is \textbf{lossy}: multiple 
distinct hash values map to the same remainder (e.g., both 
$10 \mod 8 = 2$ and $18 \mod 8 = 2$), so knowing 
$H(x) \mod m$ does not reveal what $H(x) \mod m'$ would be for 
a different table size $m'$. This lost information is precisely what 
prevents cuckoo filters from expanding. In a quotient filter, the quotient and remainder reconstruct the original value exactly, so no 
information is ever lost.

% ---- EXAMPLE: Hash split diagram ----
\begin{figure}[H]
  \centering
  \begin{tikzpicture}
    % Full hash
    \node[font=\small\ttfamily] at (-1.5, 0) {$h(x) =$};
    
    % Quotient bits
    \foreach \i/\val in {0/1, 1/0, 2/1} {
      \node[bit, fill=AccentTeal!20] (q\i) at (\i*0.6, 0) {\val};
    }
    
    % Remainder bits
    \foreach \i/\val in {3/1, 4/0, 5/1, 6/1, 7/0} {
      \node[bit, fill=AccentOrange!20] (r\i) at (\i*0.6, 0) {\val};
    }
    
    % Braces
    \draw[brace] (q0.north west) -- (q2.north east) 
      node[midway, above=10pt, font=\small, text=AccentTeal] 
      {Quotient $h_q(x)$ ($q = 3$ bits)};
    \draw[brace] (r3.north west) -- (r7.north east) 
      node[midway, above=10pt, yshift=20px, font=\small, text=AccentOrange] 
      {Remainder $h_r(x)$ ($r = 5$ bits)};
    
    % Annotations
    \node[lbl, align=center] at (1*0.6, -1.1) 
      {$\rightarrow$ slot address\\$\quot{101}_2 = 5$};
    \node[lbl, align=center, xshift=20px] at (5*0.6, -1.1) 
      {$\rightarrow$ stored fingerprint\\$\rem{10110}_2$};
  \end{tikzpicture}
  \caption{An 8-bit hash split into a 3-bit quotient (slot address) and 
           5-bit remainder (stored fingerprint). The table has 
           $2^3 = 8$ slots.}
  \label{fig:quotienting}
\end{figure}

\begin{insightbox}[Cuckoo vs.\ Quotient: The Key Difference]
In a \textbf{cuckoo filter}, the bucket index and fingerprint are 
\emph{independent}, which means knowing one tells you nothing about the other. 
In a \textbf{quotient filter}, they are \emph{complementary}, and 
together they reconstruct the full hash. This complementarity is what 
enables expansion without access to the original keys.
\end{insightbox}

\subsection{Worked Example: Insertion}

Consider a quotient filter with $p = 8$-bit hashes, $q = 3$-bit 
quotients, and $r = 5$-bit remainders. The table has $2^3 = 8$ slots. 
Suppose we insert three elements:

\begin{align*}
  h(\texttt{"hello"}) &= \underbrace{101}_{\text{quotient}} 
                         \underbrace{10101}_{\text{remainder}} 
                         \quad \rightarrow \text{slot } 5, 
                         \text{ store } \rem{10101} \\
  h(\texttt{"world"}) &= \underbrace{011}_{\text{quotient}} 
                         \underbrace{00110}_{\text{remainder}} 
                         \quad \rightarrow \text{slot } 3, 
                         \text{ store } \rem{00110} \\
  h(\texttt{"cat"})   &= \underbrace{101}_{\text{quotient}} 
                         \underbrace{11010}_{\text{remainder}} 
                         \quad \rightarrow \text{slot } 5, 
                         \text{ store } \rem{11010}
\end{align*}

The first two elements go directly into their home slots. But 
\texttt{"cat"} also hashes to slot 5, which is already occupied by 
\texttt{"hello"}. This is a \textbf{soft collision}: different keys with 
the same quotient but different remainders. The quotient filter resolves 
this by storing \texttt{"cat"}'s remainder in the next available slot 
(slot 6) using linear probing. Remainders within a run are kept in 
sorted order, so \texttt{"hello"}'s remainder (10101) stays in the 
home slot and \texttt{"cat"}'s larger remainder (11010) is placed after it:

% ---- Diagram: Simple insertion with collision ----
\begin{figure}[H]
  \centering
  \begin{tikzpicture}
    % Slots
    \foreach \i in {0,...,7} {
      \node[emptyslot, minimum width=1.3cm] (s\i) at (\i*1.4, 0) {};
      \node[font=\tiny, text=MedGray] at (\i*1.4, -0.55) {\i};
    }
    
    % Filled slots
    \node[fullslot, minimum width=1.3cm] at (3*1.4, 0) {\rem{00110}};
    \node[fullslot, minimum width=1.3cm] at (5*1.4, 0) {\rem{10101}};
    \node[fullslot, minimum width=1.3cm] at (6*1.4, 0) {\rem{11010}};
    
    % Labels
    \node[font=\scriptsize, text=AccentTeal] at (3*1.4, 0.6) 
      {\texttt{"world"}};
    \node[font=\scriptsize, text=AccentTeal] at (5*1.4, 0.6) 
      {\texttt{"hello"}};
    \node[font=\scriptsize, text=AccentTeal] at (6*1.4, 0.6) 
      {\texttt{"cat"}};
    
    % Show that cat was shifted
    \draw[arrow, AccentRed, thick] (5*1.4 - 0.1, -0.9) -- (6*1.4 + 0.1, -0.9)
      node[midway, below, font=\scriptsize, text=AccentRed, yshift=-3pt] 
      {shifted right};
    
    % Brace for run
    \draw[brace] (5*1.4 - 0.6, 1.0) -- (6*1.4 + 0.6, 1.0)
      node[midway, above=10pt, font=\scriptsize, text=DarkGray] 
      {run for quotient 5};
  \end{tikzpicture}
  \caption{After inserting \texttt{"hello"}, \texttt{"world"}, and 
           \texttt{"cat"}. Elements with the same quotient form a 
           \textbf{run} stored in contiguous slots.}
  \label{fig:quotient-insertion}
\end{figure}

The set of remainders sharing the same quotient is called a \textbf{run}. 
A sequence of contiguous runs is called a \textbf{cluster}; as the table 
fills, clusters grow longer, which increases lookup time, a limitation 
we will revisit.


\section{Metadata Bits}

Storing remainders in shifted positions creates an ambiguity: given a 
remainder in some slot, how do we know which quotient it belongs to? The 
quotient filter resolves this with three metadata bits per slot 
\cite{quotient}:

\begin{definitionbox}[Metadata Bits]
Each slot $j$ in the quotient filter maintains three single-bit flags:
\begin{enumerate}
  \item \textbf{\texttt{is\_occupied}}: Set to 1 if quotient $j$ is the 
        home slot for \emph{some} element in the filter (the element may 
        be stored elsewhere due to shifting).
  \item \textbf{\texttt{is\_continuation}}: Set to 1 if this slot holds 
        a remainder that is \emph{not the first} in its run (i.e., it 
        continues a run from a previous slot).
  \item \textbf{\texttt{is\_shifted}}: Set to 1 if the remainder in this 
        slot has been displaced from its home slot (i.e., it does not 
        occupy the slot matching its quotient).
\end{enumerate}
\end{definitionbox}

These three bits encode enough information to reconstruct which quotient 
each remainder belongs to, without storing the quotient explicitly. The 
key relationships are:

\begin{itemize}
  \item A slot with \texttt{is\_occupied = 0}, \texttt{is\_continuation = 0}, 
        \texttt{is\_shifted = 0} is \textbf{empty}.
  \item A slot with \texttt{is\_shifted = 0} and 
        \texttt{is\_continuation = 0} marks the \textbf{start of a 
        cluster}, meaning it is the first element of a run in its home slot.
  \item \texttt{is\_continuation = 0} marks the \textbf{start of a new 
        run} within a cluster.
  \item \texttt{is\_occupied} is attached to the \emph{slot index}, not 
        the stored remainder. It says ``some element with this quotient 
        exists somewhere in the table.''
\end{itemize}

\subsection{Worked Example: Metadata in Action}

Continuing from Figure~\ref{fig:quotient-insertion}, suppose we also 
insert an element with quotient 6:

\begin{align*}
  h(\texttt{"dog"}) &= \underbrace{110}_{\text{quotient}} 
                       \underbrace{01001}_{\text{remainder}} 
                       \quad \rightarrow \text{slot } 6
\end{align*}

Slot 6 is occupied by \texttt{"cat"} (which was shifted from slot 5). 
So \texttt{"dog"}'s remainder gets pushed to slot 7. The full state:

% ---- Diagram: Metadata bits ----
\begin{figure}[H]
  \centering
  \begin{tikzpicture}
    % Column headers
    \node[font=\scriptsize\bfseries, text=DarkGray] at (-2.2, 0) {Slot};
    \node[font=\scriptsize\bfseries, text=DarkGray] at (-2.2, -0.7) {Remainder};
    \node[font=\scriptsize\bfseries, text=DarkGray] at (-2.2, -1.4) 
      {\texttt{is\_occupied}};
    \node[font=\scriptsize\bfseries, text=DarkGray] at (-2.2, -2.1) 
      {\texttt{is\_cont.}};
    \node[font=\scriptsize\bfseries, text=DarkGray] at (-2.2, -2.8) 
      {\texttt{is\_shifted}};
    \node[font=\scriptsize\bfseries, text=DarkGray] at (-2.2, -3.6) 
      {Belongs to};
    
    % Slot 3 — "world"
    \node[font=\scriptsize] at (0, 0) {3};
    \node[fullslot, minimum width=1.3cm] at (0, -0.7) {\rem{00110}};
    \node[font=\scriptsize] at (0, -1.4) {1};
    \node[font=\scriptsize] at (0, -2.1) {0};
    \node[font=\scriptsize] at (0, -2.8) {0};
    \node[font=\scriptsize, text=AccentTeal] at (0, -3.6) {$q = 3$};
    
    % Slot 4 — empty
    \node[font=\scriptsize] at (1.6, 0) {4};
    \node[emptyslot, minimum width=1.3cm] at (1.6, -0.7) {};
    \node[font=\scriptsize, text=MedGray] at (1.6, -1.4) {0};
    \node[font=\scriptsize, text=MedGray] at (1.6, -2.1) {0};
    \node[font=\scriptsize, text=MedGray] at (1.6, -2.8) {0};
    \node[font=\scriptsize, text=MedGray] at (1.6, -3.6) {---};
    
    % Slot 5 — "hello"
    \node[font=\scriptsize] at (3.2, 0) {5};
    \node[fullslot, minimum width=1.3cm] at (3.2, -0.7) {\rem{10101}};
    \node[font=\scriptsize] at (3.2, -1.4) {1};
    \node[font=\scriptsize] at (3.2, -2.1) {0};
    \node[font=\scriptsize] at (3.2, -2.8) {0};
    \node[font=\scriptsize, text=AccentTeal] at (3.2, -3.6) {$q = 5$};
    
    % Slot 6 — "cat" (shifted from 5)
    \node[font=\scriptsize] at (4.8, 0) {6};
    \node[fullslot, minimum width=1.3cm] at (4.8, -0.7) {\rem{11010}};
    \node[font=\scriptsize] at (4.8, -1.4) {1};
    \node[font=\scriptsize] at (4.8, -2.1) {1};
    \node[font=\scriptsize] at (4.8, -2.8) {1};
    \node[font=\scriptsize, text=AccentTeal] at (4.8, -3.6) {$q = 5$};
    
    % Slot 7 — "dog" (shifted from 6)
    \node[font=\scriptsize] at (6.4, 0) {7};
    \node[fullslot, minimum width=1.3cm] at (6.4, -0.7) {\rem{01001}};
    \node[font=\scriptsize] at (6.4, -1.4) {0};
    \node[font=\scriptsize] at (6.4, -2.1) {0};
    \node[font=\scriptsize] at (6.4, -2.8) {1};
    \node[font=\scriptsize, text=AccentTeal] at (6.4, -3.6) {$q = 6$};
    
    % Run braces
    \draw[brace] (3.2 - 0.6, 0.4) -- (4.8 + 0.6, 0.4)
      node[midway, above=10pt, font=\scriptsize, text=AccentBlue] 
      {run for $q = 5$};
    \draw[decorate, decoration={brace, amplitude=4pt, raise=2pt, 
      mirror}, thick, DarkGray] 
      (3.2 - 0.6, -4.1) -- (6.4 + 0.6, -4.1)
      node[midway, below=8pt, font=\scriptsize, text=DarkGray] 
      {cluster (contiguous occupied region)};
  \end{tikzpicture}
  \caption{Quotient filter state after inserting \texttt{"world"} 
           ($q = 3$), \texttt{"hello"} ($q = 5$), \texttt{"cat"} 
           ($q = 5$), and \texttt{"dog"} ($q = 6$). The metadata bits 
           encode run boundaries and shift status.}
  \label{fig:quotient-metadata}
\end{figure}

Reading the metadata for each slot:
\begin{itemize}
  \item \textbf{Slot 3}: \texttt{is\_occupied = 1} (quotient 3 has 
        elements), \texttt{is\_shifted = 0} (it is in its home slot), 
        \texttt{is\_continuation = 0} (first in its run). This is 
        \texttt{"world"}'s home.
  \item \textbf{Slot 5}: \texttt{is\_occupied = 1} (quotient 5 has 
        elements), \texttt{is\_shifted = 0} (in its home slot), 
        \texttt{is\_continuation = 0} (starts a new run). This is 
        \texttt{"hello"}.
  \item \textbf{Slot 6}: \texttt{is\_occupied = 1} (quotient 6 has 
        elements somewhere), \texttt{is\_shifted = 1} (not in its home 
        slot), \texttt{is\_continuation = 1} (continues the run from 
        slot 5). This is \texttt{"cat"}, which belongs to quotient 5 but 
        is stored in slot 6.
  \item \textbf{Slot 7}: \texttt{is\_occupied = 0} (no element has 
        quotient 7), \texttt{is\_shifted = 1} (displaced from slot 6), 
        \texttt{is\_continuation = 0} (starts a new run). This is 
        \texttt{"dog"}, which belongs to quotient 6.
\end{itemize}

\subsection{Lookup Algorithm}

To query whether an element $x$ with quotient $q$ and remainder $r$ is 
in the filter:

\begin{enumerate}
  \item Check slot $q$. If \texttt{is\_occupied = 0}, then no element 
        with quotient $q$ exists, return \textsc{no}.
  \item \textbf{Find the cluster start}: scan left from slot $q$ until 
        a slot with \texttt{is\_shifted = 0} is found. This is where 
        the cluster begins.
  \item \textbf{Skip preceding runs}: scan right through the cluster, 
        counting runs. Each \texttt{is\_occupied = 1} slot to the left 
        of $q$ indicates a run to skip. Each 
        \texttt{is\_continuation = 0} slot marks the start of a new run. 
        When the count reaches zero, the current position is the start 
        of quotient $q$'s run.
  \item \textbf{Search the run}: compare the remainder in each slot of 
        the run to $r$. If a match is found, return \textsc{yes} 
        (probably in the set). If the run ends without a match, return 
        \textsc{no} (definitely not in the set).
\end{enumerate}

\begin{examplebox}[Lookup for \texttt{"dog"} (quotient $= 6$)]
\begin{enumerate}
  \item Slot 6: \texttt{is\_occupied = 1} $\rightarrow$ quotient 6 has 
        elements. Proceed.
  \item Scan left: slot 6 has \texttt{is\_shifted = 1}, slot 5 has 
        \texttt{is\_shifted = 0} $\rightarrow$ cluster starts at slot 5.
  \item Scan right from slot 5, counting runs. Between slot 5 and slot 6, 
        there is one \texttt{is\_occupied} slot (slot 5) before our 
        target (slot 6), so we skip one run:
        \begin{itemize}
          \item Slot 5: \texttt{is\_continuation = 0} $\rightarrow$ 
                start of run for $q = 5$ (skip this run)
          \item Slot 6: \texttt{is\_continuation = 1} $\rightarrow$ 
                still $q = 5$'s run
          \item Slot 7: \texttt{is\_continuation = 0} $\rightarrow$ 
                new run starts $\rightarrow$ this is $q = 6$'s run!
        \end{itemize}
  \item Slot 7 contains remainder $\rem{01001}$. Compare to 
        \texttt{"dog"}'s remainder $\rem{01001}$ $\rightarrow$ match! 
        Return \textsc{yes}.
\end{enumerate}
\end{examplebox}

Notice that this lookup required scanning through multiple slots. In the 
worst case, if all elements hash to the same quotient, a single run 
could span the entire table, requiring $O(n)$ time. In practice with a 
good hash function, runs are short and expected lookup time is $O(1)$. 
However, as the load factor increases, clusters grow and performance 
degrades; quotient filters are typically not filled beyond 75\% 
\cite{quotient}.


\section{How Expansion Works}

This is where the quotient filter's design pays off. Because the 
quotient and remainder together form the complete hash, expanding the 
table requires no access to the original keys, only a reinterpretation 
of the stored information.

\subsection{The Expansion Procedure}

To double the table from $2^q$ to $2^{q+1}$ slots:

\begin{enumerate}
  \item \textbf{Move the split point}: the quotient grows from $q$ bits 
        to $q + 1$ bits, and the remainder shrinks from $r$ bits to 
        $r - 1$ bits.
  \item \textbf{Redistribute entries}: for each stored remainder, the 
        topmost bit of the old remainder becomes part of the new quotient. 
        If this bit is 0, the entry stays in its current slot. If it is 
        1, the entry moves to the corresponding slot in the new (upper) 
        half of the table.
  \item \textbf{Truncate remainders}: remove the top bit from each 
        remainder, since it is now part of the quotient.
\end{enumerate}

\begin{examplebox}[Expansion Step by Step]
Before expansion ($q = 3$, $r = 5$, table has $2^3 = 8$ slots):
\[
  h(\texttt{"hello"}) = \underbrace{101}_{\text{quotient}} 
                         \underbrace{10101}_{\text{remainder}}
  \quad \rightarrow \text{slot } 5, \text{ store } \rem{10101}
\]

After expansion ($q = 4$, $r = 4$, table has $2^4 = 16$ slots):
\[
  h(\texttt{"hello"}) = \underbrace{1011}_{\text{new quotient}} 
                         \underbrace{0101}_{\text{new remainder}}
  \quad \rightarrow \text{slot } 11, \text{ store } \rem{0101}
\]

The hash \texttt{10110101} did not change. We simply read one more bit 
as part of the quotient. The top bit of the old remainder (\texttt{1}) 
became the new lowest bit of the quotient, changing the slot from 
$\texttt{101}_2 = 5$ to $\texttt{1011}_2 = 11$.
\end{examplebox}

% ---- Diagram: Expansion ----
\begin{figure}[H]
  \centering
  \begin{tikzpicture}
    % BEFORE: show the hash split
    \node[lbl, font=\bfseries] at (-2.5, 2.5) {Before ($q = 3$, $r = 5$):};
    
    \node[font=\small\ttfamily] at (-1.5, 1.7) {$h(x) =$};
    \foreach \i/\val in {0/1, 1/0, 2/1} {
      \node[bit, fill=AccentTeal!20] at (\i*0.5, 1.7) {\val};
    }
    \foreach \i/\val in {3/1, 4/0, 5/1, 6/0, 7/1} {
      \node[bit, fill=AccentOrange!20] at (\i*0.5, 1.7) {\val};
    }
    \draw[brace] (0*0.5-0.2, 2.15) -- (2*0.5+0.2, 2.15)
      node[midway, above=10pt, font=\scriptsize, text=AccentTeal] {$q$: slot 5};
    \draw[brace] (3*0.5-0.2, 2.15) -- (7*0.5+0.2, 2.15)
      node[midway, above=10pt, font=\scriptsize, text=AccentOrange] 
      {$r$: \rem{10101}};
    
    % Arrow
    \draw[arrow, thick, AccentBlue] (2, 1.0) -- (2, 0.3)
  node[midway, right=4pt, font=\small, align=left] {steals 1 bit\\from r to q};

    
    % AFTER: new split
    \node[lbl, font=\bfseries] at (-2.5, -0.2) {After ($q = 4$, $r = 4$):};
    
    \node[font=\small\ttfamily] at (-1.5, -1.0) {$h(x) =$};
    \foreach \i/\val in {0/1, 1/0, 2/1, 3/1} {
      \node[bit, fill=AccentTeal!20] at (\i*0.5, -1.0) {\val};
    }
    \foreach \i/\val in {4/0, 5/1, 6/0, 7/1} {
      \node[bit, fill=AccentOrange!20] at (\i*0.5, -1.0) {\val};
    }
    \draw[brace] (0*0.5-0.2, -0.55) -- (3*0.5+0.2, -0.55)
      node[midway, above=10pt, font=\scriptsize, text=AccentTeal] 
      {$q$: slot 11};
    \draw[brace] (4*0.5-0.2, -0.55) -- (7*0.5+0.2, -0.55)
      node[midway, above=10pt, font=\scriptsize, text=AccentOrange] 
      {$r$: \rem{0101}};
    
    % Highlight the stolen bit
    \draw[AccentRed, very thick, rounded corners=2pt] 
      (3*0.5-0.22, 1.7-0.28) rectangle (3*0.5+0.22, 1.7+0.28);
    \draw[AccentRed, very thick, rounded corners=2pt] 
      (3*0.5-0.22, -1.0-0.28) rectangle (3*0.5+0.22, -1.0+0.28);
    \draw[arrow, AccentRed, dashed] (3*0.5+0.3, 1.42) -- (3*0.5+0.3, -0.72);
    
    % Summary
    \node[lbl, align=left, font=\small] at (6.5, 0.5) {
      Same hash: \texttt{10110101}\\
      Slot: $5 \rightarrow 11$\\
      Remainder: 5 bits $\rightarrow$ 4 bits\\
      No original key needed!
    };
  \end{tikzpicture}
  \caption{Expansion moves the split point one bit to the right. The top 
           bit of the old remainder (highlighted) becomes the new lowest 
           bit of the quotient. The hash itself is unchanged.}
  \label{fig:expansion-split}
\end{figure}

% ---- Diagram: Table before and after ----
\begin{figure}[H]
  \centering
  \begin{tikzpicture}[node distance=0pt]
    % BEFORE expansion
    \node[lbl, font=\bfseries] at (1.5, 1.8) {Before Expansion (8 slots)};
    
    \foreach \i/\fp in {0/, 1/, 2/, 3/\rem{00110}, 4/, 
                        5/\rem{10101}, 6/\rem{11010}, 7/} {
      \ifx\fp\empty
        \node[emptyslot, minimum width=1.3cm] (a\i) at (\i*1.4, 1) {};
      \else
        \node[fullslot, minimum width=1.3cm] (a\i) at (\i*1.4, 1) {\fp};
      \fi
      \node[font=\tiny, text=MedGray] at (\i*1.4, 0.45) {\i};
    }
    
    % Arrow down
\begin{scope}[yshift=-10pt]
  \draw[arrow, thick, AccentBlue] (4.9, 0.5) -- (4.9, -0.5) 
    node[midway, right, font=\small] 
    {expand: steal 1 bit from each remainder};
\end{scope}    
    % AFTER expansion (16 slots, show only relevant ones)
  \begin{scope}[yshift=-10pt]
     \node[lbl, font=\bfseries] at (4.5, -1.0) {After Expansion (16 slots)};
    
    \foreach \i/\fp in {0/, 1/, 2/, 3/, 4/, 5/, 6/\rem{0110}, 
                        7/, 8/, 9/, 10/, 11/\rem{0101}, 
                        12/, 13/\rem{1010}} {
      \ifx\fp\empty
        \node[emptyslot, minimum width=0.75cm] (b\i) at (\i*0.8, -1.8) {};
      \else
        \node[fullslot, minimum width=0.75cm] (b\i) at (\i*0.8, -1.8) {\fp};
      \fi
      \node[font=\tiny, text=MedGray] at (\i*0.8, -2.35) {\i};
    }
    
    \node[lbl, align=center] at (4.5, -3.1) {
      Remainders shrink: 5 bits $\rightarrow$ 4 bits.
      Entries redistribute across doubled table.
    };
  \end{scope}
  \end{tikzpicture}
  \caption{Quotient filter expansion: the table doubles from 8 to 16 
           slots. Each remainder loses its top bit to the quotient.
           \texttt{"hello"} moves from slot 5 to slot 11; 
           \texttt{"cat"} moves from slot 6 to slot 13.}
  \label{fig:expansion}
\end{figure}

\begin{insightbox}[Why This Works Without Original Keys]
The expansion procedure never needs the original element $x$ or the 
full hash $\text{hash}(x)$. It only needs the \emph{quotient} (known 
from the slot position) and the \emph{remainder} (stored in the slot). 
Together these reconstruct the full $p$-bit hash, which can then be 
re-split at the new position. This is the fundamental advantage of 
quotienting over the cuckoo filter's modulo-based addressing.
\end{insightbox}


\section{The Void Entry Problem}

Expansion comes at a cost: each time the table doubles, every remainder 
loses one bit. For an entry that was present at the very first expansion 
and survives through $X$ expansions, its remainder shrinks from $r$ bits 
to $r - X$ bits.

\subsection{How Void Entries Arise}

Consider an entry inserted with an initial remainder of $r = 5$ bits. 
After each expansion, the top bit of its remainder is consumed by the 
growing quotient:

% ---- Diagram: Void entry formation ----
\begin{figure}[H]
  \centering
  \begin{tikzpicture}
    \node[lbl, font=\bfseries] at (-2.5, 1) {Expansion};
    \node[lbl, font=\bfseries] at (1.5, 1) {Remainder};
    \node[lbl, font=\bfseries] at (4.5, 1) {Bits left};
    
    % Row 0
    \node[lbl] at (-2.5, 0) {0 (initial)};
    \node[bit, fill=BitOne] at (0.5, 0) {1};
    \node[bit, fill=BitOne] at (1.0, 0) {0};
    \node[bit, fill=BitOne] at (1.5, 0) {1};
    \node[bit, fill=BitOne] at (2.0, 0) {0};
    \node[bit, fill=BitOne] at (2.5, 0) {1};
    \node[lbl] at (4.5, 0) {5};
    
    % Row 1
    \node[lbl] at (-2.5, -0.8) {1};
    \node[bit, fill=BitOne] at (1.0, -0.8) {0};
    \node[bit, fill=BitOne] at (1.5, -0.8) {1};
    \node[bit, fill=BitOne] at (2.0, -0.8) {0};
    \node[bit, fill=BitOne] at (2.5, -0.8) {1};
    \node[bit, fill=BitZero, opacity=0.3] at (0.5, -0.8) {};
    \node[lbl] at (4.5, -0.8) {4};
    
    % Row 2
    \node[lbl] at (-2.5, -1.6) {2};
    \node[bit, fill=BitOne] at (1.5, -1.6) {1};
    \node[bit, fill=BitOne] at (2.0, -1.6) {0};
    \node[bit, fill=BitOne] at (2.5, -1.6) {1};
    \node[bit, fill=BitZero, opacity=0.3] at (0.5, -1.6) {};
    \node[bit, fill=BitZero, opacity=0.3] at (1.0, -1.6) {};
    \node[lbl] at (4.5, -1.6) {3};
    
    % Row 3
    \node[lbl] at (-2.5, -2.4) {3};
    \node[bit, fill=BitOne] at (2.0, -2.4) {0};
    \node[bit, fill=BitOne] at (2.5, -2.4) {1};
    \node[bit, fill=BitZero, opacity=0.3] at (0.5, -2.4) {};
    \node[bit, fill=BitZero, opacity=0.3] at (1.0, -2.4) {};
    \node[bit, fill=BitZero, opacity=0.3] at (1.5, -2.4) {};
    \node[lbl] at (4.5, -2.4) {2};
    
    % Row 4
    \node[lbl] at (-2.5, -3.2) {4};
    \node[bit, fill=BitOne] at (2.5, -3.2) {1};
    \node[bit, fill=BitZero, opacity=0.3] at (0.5, -3.2) {};
    \node[bit, fill=BitZero, opacity=0.3] at (1.0, -3.2) {};
    \node[bit, fill=BitZero, opacity=0.3] at (1.5, -3.2) {};
    \node[bit, fill=BitZero, opacity=0.3] at (2.0, -3.2) {};
    \node[lbl] at (4.5, -3.2) {1};
    
    % Row 5 — VOID
    \node[lbl, text=AccentRed, font=\bfseries] at (-2.5, -4.0) {5};
    \node[bit, fill=VoidColor] at (0.5, -4.0) {?};
    \node[bit, fill=VoidColor] at (1.0, -4.0) {?};
    \node[bit, fill=VoidColor] at (1.5, -4.0) {?};
    \node[bit, fill=VoidColor] at (2.0, -4.0) {?};
    \node[bit, fill=VoidColor] at (2.5, -4.0) {?};
    \node[lbl, text=AccentRed, font=\bfseries] at (4.5, -4.0) {0 = \void{}};
    
    % Arrow
    \draw[arrow, AccentRed, very thick] (5.5, -4.0) -- (7.5, -4.0) 
      node[right, font=\small, text=AccentRed, align=left] {
        No fingerprint bits left!\\
        Matches \textbf{any} query\\
        $\rightarrow$ 100\% false positive};
  \end{tikzpicture}
  \caption{A remainder shrinks by one bit per expansion. After $r$ 
           expansions, zero bits remain, the entry becomes \void{}.}
  \label{fig:void-formation}
\end{figure}

\begin{definitionbox}[Void Entry]
A \textbf{void entry} is an entry whose remainder has been reduced to 
zero bits through repeated expansions. Since there is no fingerprint 
left to compare against, a void entry matches \emph{every} query to 
its slot, producing a guaranteed false positive. Formally, if an entry 
was created with remainder length $r$ and has survived $X \geq r$ 
expansions, its current remainder length is $\max(0, r - X) = 0$.
\end{definitionbox}

\subsection{The Expansion Dilemma}

Void entries create a specific problem during expansion. When the table 
doubles, each entry moves to one of two possible slots in the new table, determined by the top bit of its current remainder. But a void entry 
has \emph{no remaining bits}. Without a bit to read, the filter cannot 
determine which of the two new slots the entry belongs in:

% ---- Diagram: Void entry expansion problem ----
\begin{figure}[H]
  \centering
  \begin{tikzpicture}
    % Before
    \node[lbl, font=\bfseries] at (1.5, 2) {Before expansion (8 slots)};
    \foreach \i in {0,...,7} {
      \node[emptyslot, minimum width=0.9cm] (a\i) at (\i*1.0, 1.2) {};
      \node[font=\tiny, text=MedGray] at (\i*1.0, 0.75) {\i};
    }
    \node[voidslot, minimum width=0.9cm] at (3*1.0, 1.2) {V};
    
    % After
    \node[lbl, font=\bfseries] at (5, -2.7) {After expansion (16 slots)};
    \foreach \i in {0,...,15} {
      \node[emptyslot, minimum width=0.55cm] (b\i) at (\i*0.6, -1.5) {};
      \node[font=\tiny, text=MedGray] at (\i*0.6, -1.95) {\i};
    }
    
    % Two possible destinations
    \node[voidslot, minimum width=0.55cm] at (6*0.6, -1.5) {?};
    \node[voidslot, minimum width=0.55cm] at (7*0.6, -1.5) {?};
    
    % Arrows to both
    \draw[dasharrow, AccentRed] (3*1.0, 0.85) -- (6*0.6, -1.15)
      node[midway, left, xshift=-4pt, font=\scriptsize, text=AccentRed] {bit = 0?};
    \draw[dasharrow, AccentRed] (3*1.0, 0.85) -- (7*0.6, -1.15)
      node[midway, right, xshift=4pt, font=\scriptsize, text=AccentRed] {bit = 1?};
    
    % Annotation
    \node[lbl, text=AccentRed, font=\small, align=center] at (12, -0.1) 
      {No remainder bits left\\to decide which slot!\\\\
       Slot 6 (if bit was 0)\\or slot 7 (if bit was 1)\\
       \textbf{We don't know.}};
  \end{tikzpicture}
  \caption{A void entry at slot 3 must move to either slot 6 or slot 7 
           during expansion, but with zero remainder bits, there is no 
           information to determine which. This is the core problem that 
           InfiniFilter and the Aleph Filter each solve differently.}
  \label{fig:void-expansion}
\end{figure}

\subsection{InfiniFilter's Approach: Secondary Tables}

InfiniFilter \cite{infinifilter} solves this by placing void entries into 
\textbf{secondary hash tables}, one per expansion level. When an entry 
becomes void, it is moved to a secondary table that stores enough 
additional hash bits to track its position.

However, this means that queries for non-existent or old keys must 
now search the main table \emph{and} all secondary tables:

% ---- Diagram: InfiniFilter multi-table query ----
\begin{figure}[H]
  \centering
  \begin{tikzpicture}
    % Main table
    \node[lbl, font=\bfseries] at (-1.5, 0) {Main table};
    \foreach \i in {0,...,7} {
      \node[emptyslot, minimum width=0.7cm] (m\i) at (\i*0.8, -0.7) {};
    }
    
    % Secondary table 1
    \node[lbl, font=\bfseries] at (-1.5, -1.8) {Secondary 1};
    \foreach \i in {0,...,3} {
      \node[emptyslot, minimum width=0.7cm, fill=orange!8] 
        (s\i) at (\i*0.8, -2.5) {};
    }
    
    % Secondary table 2
    \node[lbl, font=\bfseries] at (-1.5, -3.6) {Secondary 2};
    \foreach \i in {0,...,1} {
      \node[emptyslot, minimum width=0.7cm, fill=red!8] 
        (t\i) at (\i*0.8, -4.3) {};
    }
    
    % Query path
    \node[hashbox] (q) at (8, 0) {query $x$};
    \draw[arrow, AccentBlue] (q.west) -- (m7.east)
      node[midway, above, yshift=10pt, font=\scriptsize] {check};
    \draw[dasharrow, AccentRed] (m7.south) -- (s3.north east)
      node[midway, yshift=-6pt, right, font=\scriptsize, text=AccentRed] {miss};
    \draw[dasharrow, AccentRed] (s3.south) -- (t1.north east)
      node[midway, yshift=-6pt, right, font=\scriptsize, text=AccentRed] {miss};
    
    % Annotation
    \node[lbl, text=AccentRed, align=left, font=\small] at (7.5, -3.5) {
      Must check \textbf{every} table\\
      Query time: $\bigO{\log(N)/F}$\\
      (one table per expansion)
    };
  \end{tikzpicture}
  \caption{InfiniFilter stores void entries in secondary hash tables. 
           In the fixed-width regime, queries may need to search all tables,
           degrading to $\bigO{\log(N)/F}$ (equivalently $\bigO{\log N}$ for
           constant $F$).}
  \label{fig:infinifilter}
\end{figure}

With each expansion adding a new secondary table, the number of tables 
grows as $\bigO{\log(N)/F}$ in the fixed-width regime (where $N$ is the 
total number of insertions and $F$ is the base fingerprint length). 
A query that targets a non-existent key or an old key must search every 
table, giving worst-case $\bigO{\log(N)/F}$ query time (or $\bigO{\log N}$ 
for constant $F$), a significant 
degradation from the $\bigO{1}$ queries of the original quotient filter.

\subsection{The Aleph Filter's Solution: Void Duplication}

The Aleph Filter \cite{aleph} takes a fundamentally different approach. 
Instead of moving void entries to secondary tables, it \textbf{duplicates} 
each void entry into \emph{both} possible slots during expansion. The two candidate slots in the expanded table are $2q$ and $2q + 1$, 
corresponding to appending a $0$ or $1$ bit to the old quotient:
\[
  q \| 0 = 2q \qquad \text{and} \qquad q \| 1 = 2q + 1
\]
where $\|$ denotes bit concatenation. A comparison betwen InfiniFilter and Aleph Filter:

% ---- Diagram: Aleph vs InfiniFilter ----
\begin{figure}[H]
  \centering
  \begin{tikzpicture}[scale=0.9, every node/.style={transform shape}]
    % BEFORE
    \node[lbl, font=\bfseries] at (1.6, 1.5) {Before Expansion};
    \foreach \i/\content/\style in {0//emptyslot, 1//emptyslot, 
                                     2//emptyslot, 3/V/voidslot} {
      \node[\style, minimum width=1cm] (a\i) at (\i*1.2, 0.7) {\content};
      \node[font=\tiny, text=MedGray] at (\i*1.2, 0.15) {\i};
    }
    
    % Arrow
    \draw[arrow, very thick, AccentBlue] (2, -0.3) -- (2, -1.0);
    
    % AFTER — InfiniFilter
    \node[lbl, font=\bfseries] at (-3, -1.6) {InfiniFilter:};
    \foreach \i in {0,...,7} {
      \node[emptyslot, minimum width=0.8cm] (inf\i) at ({-5+\i*0.9}, -2.3) {};
    }
    \node[font=\tiny, text=MedGray] at (-5, -2.8) {0};
    \node[font=\tiny, text=MedGray] at ({-5+7*0.9}, -2.8) {7};
    
    \node[voidslot, minimum width=0.8cm] (sec) at (-3, -3.5) {V};
    \node[lbl, font=\scriptsize] at (-3, -4.1) {secondary table};
    \draw[dasharrow, AccentRed] (inf3.south) -- (sec.north);
    \node[lbl, text=AccentRed, font=\scriptsize, anchor=west] at (0, -3.5) 
      {Query: $\bigO{\log N}$};
    
    % AFTER — Aleph
    \node[lbl, font=\bfseries] at (5.5, -1.6) {Aleph Filter:};
    \foreach \i in {0,...,7} {
      \node[emptyslot, minimum width=0.8cm] (al\i) at ({3+\i*0.9}, -2.3) {};
    }
    \node[font=\tiny, text=MedGray] at (3, -2.8) {0};
    \node[font=\tiny, text=MedGray] at ({3+7*0.9}, -2.8) {7};
    
    \node[voidslot, minimum width=0.8cm] at ({3+6*0.9}, -2.3) {V};
    \node[voidslot, minimum width=0.8cm] at ({3+7*0.9}, -2.3) {V};
    
    \draw[brace] ({3+6*0.9-0.35}, -1.7) -- ({3+7*0.9+0.35}, -1.7)
      node[midway, above=10pt, font=\scriptsize, text=AccentGreen] 
      {duplicated!};
    \node[lbl, text=AccentGreen, font=\scriptsize, anchor=west, xshift=10pt] at (9.5, -2.3) 
      {Query: $\bigO{1}$};
  \end{tikzpicture}
  \caption{Aleph Filter duplicates void entries to both possible slots, 
           keeping queries $\bigO{1}$.}
  \label{fig:aleph-vs-infini}
\end{figure}

Since a void entry matches any query to its slot (it has no fingerprint 
to compare), placing a copy in both candidate slots ensures correctness: 
regardless of which slot the original element ``should'' have gone to, 
a copy is there. All entries remain in a single table, so queries never 
need to search secondary structures.

This is the core innovation of the Aleph Filter, which the next chapter 
develops in full detail. The complete design addresses three challenges:
\begin{enumerate}
  \item \textbf{Void duplication:} keeps queries $\bigO{1}$ by 
        eliminating secondary tables.
  \item \textbf{Tombstone-based deletion:} handles the complication that 
        deleting a duplicated void entry must not remove both copies.
  \item \textbf{Rejuvenation:} prevents void entries from accumulating 
        indefinitely by periodically replacing them with fresh 
        fingerprints.
\end{enumerate}

\begin{warningbox}[The Complete Filter Evolution]
\begin{center}
\begin{tabular}{l c c c c}
  \hline
  & \textbf{Delete} & \textbf{Expand} & \textbf{Query} & \textbf{Stable FPR} \\
  \hline
  Bloom filter     & \texttimes & \texttimes & $\bigO{k}$       & \texttimes \\
  Cuckoo filter    & \checkmark & \texttimes & $\bigO{1}$       & \texttimes \\
  Quotient filter  & \checkmark & \checkmark & $\bigO{1}$*      & \texttimes \\
  InfiniFilter     & \checkmark & \checkmark & $\bigO{\log(N)/F}$  & \checkmark \\
  \textbf{Aleph Filter} & \checkmark & \checkmark & $\bigO{1}$ & \checkmark \\
  \hline
\end{tabular}

\smallskip
{\small *amortized; degrades at high load factors}
\end{center}
Each row solves a limitation of the row above it. The Aleph Filter is 
the first to achieve all four properties simultaneously.
\end{warningbox}


% ============================================================================
% CHAPTER 3: THE ALEPH FILTER
% ============================================================================

\chapter{The Aleph Filter}

\section{Overview}

The previous chapter established that quotient filters can expand 
without access to the original keys by moving the quotient-remainder 
split point. However, after enough expansions, old entries run out of 
fingerprint bits and become \emph{void entries}. InfiniFilter 
\cite{infinifilter} handles void entries by moving them to secondary 
hash tables, but this degrades query time to $\bigO{\log N}$ as more 
tables accumulate.

The Aleph Filter \cite{aleph} builds directly on InfiniFilter's 
architecture, using the same slot format, variable-length fingerprints, 
and unary padding (self-delimiting age encoding), but it changes how void entries are handled during 
expansion. This single change, combined with two supporting mechanisms, 
yields $\bigO{1}$ worst-case time for all operations while expanding 
infinitely.

This chapter presents the three contributions in order. We begin with 
the slot format inherited from InfiniFilter, then introduce void 
entry duplication (Section~\ref{sec:void-duplication}), 
the core idea that enables constant-time queries 
(Section~\ref{sec:constant-queries}). We then cover tombstone-based 
deletes (Section~\ref{sec:tombstone-deletes}) and rejuvenation 
(Section~\ref{sec:rejuvenation}), two mechanisms that keep the filter 
efficient under updates. Finally, we analyze the FPR 
(Section~\ref{sec:fpr-analysis}), describe the operating regimes 
(Section~\ref{sec:regimes}), and summarize the full architecture.


\begin{insightbox}[The Three Contributions]
\begin{enumerate}
  \item \textbf{Void entry duplication:} during expansion, void 
        entries are copied to \emph{both} possible new slots, keeping 
        everything in one table $\rightarrow$ $\bigO{1}$ queries.
  \item \textbf{Tombstone-based deletes:} deleting a void entry 
        replaces it with a tombstone and defers duplicate removal to 
        the next expansion $\rightarrow$ $\bigO{1}$ deletes.
  \item \textbf{Rejuvenation:} void entries are periodically replaced 
        with fresh fingerprints, preventing unbounded accumulation of 
        duplicates $\rightarrow$ stable FPR.
\end{enumerate}
\end{insightbox}


% ============================================================================
\section{Structure}
% ============================================================================

The Aleph Filter inherits InfiniFilter's slot format. Each slot in the 
main hash table stores three metadata bits followed by a data region 
containing a variable-length fingerprint and a self-delimiting unary 
code \cite{infinifilter, aleph}:

% ---- Diagram: Slot format ----

\begin{figure}[H]
  \centering
  \begin{tikzpicture}
    % Title
    \node[lbl, font=\bfseries] at (0, 4.2) 
      {Slot format (example: 12-bit slots = 3 metadata + 9 data bits)};
    
    % New entry (long fingerprint)
    \node[lbl, anchor=east] at (-3, 2.5) {New entry:};
    \node[slot, fill=MedGray!15, minimum width=0.7cm] at (-2, 2.5) 
      {\scriptsize meta};
    \foreach \i/\val in {0/1, 1/0, 2/1, 3/1, 4/0, 5/1, 6/1, 7/0} {
      \node[bit, fill=AccentTeal!15] at ({-1.3 + \i*0.45}, 2.5) 
        {\tiny \val};
    }
    \node[bit, fill=AccentOrange!20] at ({-1.3 + 8*0.45}, 2.5) {\tiny 0};
    \draw[brace] ({-1.3 - 0.18}, 2.95) -- ({-1.3 + 7*0.45 + 0.18}, 2.95)
      node[midway, above=10pt, font=\scriptsize, text=AccentTeal] 
      {fingerprint (8 bits)};
    \node[font=\scriptsize, text=AccentOrange, anchor=west] 
      at ({-1.3 + 8*0.45 + 0.3}, 2.5) {unary};

    % Old entry (short fingerprint)
    \begin{scope}[yshift=-20pt]
    \node[lbl, anchor=east] at (-3, 1.3) {Old entry:};
    \node[slot, fill=MedGray!15, minimum width=0.7cm] at (-2, 1.3) 
      {\scriptsize meta};
    \foreach \i/\val in {0/1, 1/1, 2/0} {
      \node[bit, fill=AccentTeal!15] at ({-1.3 + \i*0.45}, 1.3) 
        {\tiny \val};
    }
    \foreach \i/\val in {3/1, 4/0, 5/0, 6/0, 7/0, 8/0} {
      \node[bit, fill=AccentOrange!20] at ({-1.3 + \i*0.45}, 1.3) 
        {\tiny \val};
    }
    \draw[brace] ({-1.3 - 0.18}, 1.75) -- ({-1.3 + 2*0.45 + 0.18}, 1.75)
      node[midway, above=10pt, font=\scriptsize, text=AccentTeal] 
      {fp (3 bits)};
    \draw[brace] ({-1.3 + 3*0.45 - 0.18}, 1.75) -- 
      ({-1.3 + 8*0.45 + 0.18}, 1.75)
      node[midway, above=10pt, font=\scriptsize, text=AccentOrange] 
      {unary (6 bits)};
    \end{scope}

    % Void entry (no fingerprint)
    \begin{scope}[yshift=-40pt]
    \node[lbl, anchor=east] at (-3, 0.1) {Void entry:};
    \node[slot, fill=MedGray!15, minimum width=0.7cm] at (-2, 0.1) 
      {\scriptsize meta};
    \foreach \i/\val in {0/1, 1/1, 2/1, 3/1, 4/1, 5/1, 6/1, 7/1} {
      \node[bit, fill=VoidColor] at ({-1.3 + \i*0.45}, 0.1) 
        {\tiny \val};
    }
    \node[bit, fill=VoidColor] at ({-1.3 + 8*0.45}, 0.1) {\tiny 0};
    \draw[brace] ({-1.3 - 0.18}, 0.55) -- ({-1.3 + 8*0.45 + 0.18}, 0.55)
      node[midway, above=10pt, font=\scriptsize, text=AccentOrange] 
      {all unary (0 fingerprint bits)};
    \end{scope}
    % Tombstone
    \begin{scope}[yshift=-60pt]

    \node[lbl, anchor=east] at (-3, -1.1) {Tombstone:};
    \node[slot, fill=MedGray!15, minimum width=0.7cm] at (-2, -1.1) 
      {\scriptsize meta};
    \foreach \i in {0,...,8} {
      \node[bit, fill=TombstoneColor] at ({-1.3 + \i*0.45}, -1.1) 
        {\tiny 1};
    }
    \draw[brace] ({-1.3 - 0.18}, -0.65) -- ({-1.3 + 8*0.45 + 0.18}, -0.65)
      node[midway, above=10pt, font=\scriptsize, text=AccentRed] 
      {all 1s (special marker)};
    \end{scope}
  \end{tikzpicture}
  \caption{The four possible slot states in an Aleph Filter. The unary 
           code is self-delimiting: read 1s until a 0 is encountered, 
           then the remaining bits are the fingerprint. A void entry has 
           no fingerprint; a tombstone is distinguished by an all-1s 
           pattern.}
  \label{fig:slot-format}
\end{figure}

The unary code allows the filter to store variable-length fingerprints 
within fixed-width slots. Newer entries receive longer fingerprints 
(more information, fewer false positives), while older entries that
have survived multiple expansions have shorter fingerprints. This is 
the key mechanism inherited from InfiniFilter that allows infinite 
expansion \cite{infinifilter}.

It is important to note that only the stored fingerprint shrinks 
across expansions. The hash function always produces a fixed-length 
output (e.g., 64 bits). When a new entry is inserted into a large 
table, the hash is split at the current table size: more bits go to 
the quotient (slot address), but the fingerprint is still initialized 
to $F$ bits. The void problem only affects \emph{old} entries, 
because they were inserted when the table was smaller and only a 
small slice of the hash was saved. The bits that were discarded at 
insertion time cannot be recovered, so after enough expansions the 
stored slice is fully consumed. New entries, by contrast, always have 
access to the full hash output and simply reach deeper into it for 
the larger quotient while still reserving $F$ bits for the fingerprint:

\[
\underbrace{q \text{ bits}}_{\log_2(\text{table size})} \;\Big|\; 
\underbrace{F \text{ bits}}_{\text{fingerprint}} \;\Big|\; 
\underbrace{\text{unused bits}}_{\text{discarded}}
\]

As the table grows, $q$ increases and the unused portion shrinks, 
but $F$ stays the same for every new insertion. In theory, if the 
table ever reached $2^{64}$ slots (far beyond any practical system), 
even new entries would have no room for fingerprints.

\subsection{The Mother Hash}
\label{sec:mother-hash}

As an entry survives expansions, the quotient-remainder split point 
slides through its hash: fingerprint bits are consumed one by one to 
extend the quotient (slot address). By the time an entry becomes void, 
its slot address in the main table contains all the hash bits that 
were ever stored for it: the original quotient bits plus every 
fingerprint bit that was transferred across expansions. This 
accumulated slot address is the entry's \textbf{mother hash}.

A mother hash inserted at Generation $j$ (when the table had 
$2^j$ slots and the quotient was $j$ bits) and that had $F$ 
fingerprint bits at insertion will have $j + F$ total bits when it 
becomes void. An entry inserted later, at Generation $j'$ where 
$j' > j$, started with a longer quotient ($j'$ bits) and the same 
$F$ fingerprint bits, giving a mother hash of $j' + F$ bits. 
Therefore, \textbf{newer entries have longer mother hashes than 
older entries}.

Understanding the mother hash is essential for the delete and 
rejuvenation mechanisms described in 
Sections~\ref{sec:tombstone-deletes} and~\ref{sec:rejuvenation}. The 
mother hash tells the filter where all of a void entry's duplicates 
are located, enabling their removal.

With the slot format and mother hash concept in place, we now turn to 
the Aleph Filter's core contribution: how it handles void entries 
during expansion.


% ============================================================================
\section{Void Entry Duplication}
\label{sec:void-duplication}
% ============================================================================

This section describes the Aleph Filter's core innovation. During 
expansion, each non-void entry moves to one of two possible slots in 
the expanded table. The lowest bit of its remainder determines which 
slot: if the old quotient is $q$, the new slot is either $2q$ (if the 
bit is 0) or $2q + 1$ (if the bit is 1). After moving, that bit is 
removed from the remainder.

A void entry has no remaining bits. It cannot determine which of the 
two new slots it belongs to. InfiniFilter solves this by moving void 
entries to secondary hash tables \cite{infinifilter}. The Aleph Filter 
takes a different approach: \textbf{duplicate the void entry into both 
possible slots} \cite{aleph}.

The intuition is straightforward: since the filter does not know 
whether the missing bit would have been 0 or 1, it places the entry 
in both branches. On the next expansion, each copy faces the same 
ambiguity and is again placed in both child slots. This continues for 
every subsequent expansion, so no matter what the original key's full 
hash would have resolved to, a void entry is guaranteed to be waiting 
at the correct slot.

\begin{definitionbox}[Void Duplication]
During expansion, for each void entry at slot $q$ in the old table, 
the Aleph Filter places a void entry at \emph{both} slot $2q$ and 
slot $2q + 1$ in the expanded table. This ensures that regardless of 
which slot the original element's hash would have directed it to, a 
matching entry is present.
\end{definitionbox}

\subsection{Worked Example}

Consider a filter with 4 slots ($2^2$) containing a void entry at 
slot 3 (binary: \texttt{11}). During expansion to 8 slots ($2^3$), 
slot 3 could map to either slot 6 (\texttt{110}) or slot 7 
(\texttt{111}):

% ---- Diagram: Void duplication step by step ----
\begin{figure}[H]
  \centering
  \begin{tikzpicture}
    % BEFORE
    \node[lbl, font=\bfseries] at (1.8, 2) {Before Expansion (4 slots)};
    \foreach \i/\content/\style in {0//emptyslot, 1//emptyslot, 
                                     2//emptyslot, 3/V/voidslot} {
      \node[\style, minimum width=1.2cm] (a\i) at (\i*1.4, 1.2) 
        {\content};
      \node[font=\tiny, text=MedGray] at (\i*1.4, 0.65) {\i};
    }
    
    % Arrow
    \draw[arrow, very thick, AccentBlue] (2.1, 0.3) -- (2.1, -0.4)
      node[midway, right=4pt, font=\small] {expand};
    
    % AFTER
    \node[lbl, font=\bfseries] at (2.5, -0.9) {After Expansion (8 slots)};
    \foreach \i in {0,...,7} {
      \node[emptyslot, minimum width=0.9cm] (b\i) at (\i*1.05, -1.7) {};
      \node[font=\tiny, text=MedGray] at (\i*1.05, -2.25) {\i};
    }
    
    % Void duplicated to BOTH slots
    \node[voidslot, minimum width=0.9cm] at (6*1.05, -1.7) {V};
    \node[voidslot, minimum width=0.9cm] at (7*1.05, -1.7) {V};
    
    % Arrows from old slot to both new slots
    \draw[arrow, AccentRed, thick] (a3.south) to[out=-70, in=90] 
      (6*1.05, -1.35);
    \draw[arrow, AccentRed, thick] (a3.south) to[out=-50, in=90] 
      (7*1.05, -1.35);
    
    % Labels
    \node[font=\scriptsize, text=AccentRed] at (4.8, -0.4) 
      {$2q = 6$};
    \node[font=\scriptsize, text=AccentRed] at (8.2, -0.5) 
      {$2q+1 = 7$};
    
    % Brace
  \draw[decorate, decoration={brace, amplitude=5pt, mirror}, thick, black] 
  (6*1.05 - 0.4, -2.5) -- (7*1.05 + 0.4, -2.5)
    node[midway, below=8pt, font=\scriptsize, text=AccentGreen] 
    {copy in both!};
  \end{tikzpicture}
  \caption{A void entry at slot 3 is duplicated to slots 6 and 7 during 
           expansion. The filter cannot determine which is correct, so 
           it places a copy in both.}
  \label{fig:void-dup-example}
\end{figure}

On the next expansion (8 $\rightarrow$ 16 slots), each void copy 
duplicates again. The copy at slot 6 (\texttt{110}) goes to slots 12 
(\texttt{1100}) and 13 (\texttt{1101}). The copy at slot 7 
(\texttt{111}) goes to slots 14 (\texttt{1110}) and 15 
(\texttt{1111}). The entry now has 4 copies:

% ---- Diagram: Exponential duplication ----
\begin{figure}[H]
  \centering
  \begin{tikzpicture}
    % Level 0: 1 copy
    \node[voidslot, minimum width=1cm] (v0) at (0, 0) {V};
    \node[lbl, font=\scriptsize, anchor=east] at (-1, 0) 
      {4 slots:};
    
    % Level 1: 2 copies
    \node[voidslot, minimum width=0.8cm] (v1a) at (-1.5, -1.5) {V};
    \node[voidslot, minimum width=0.8cm] (v1b) at (1.5, -1.5) {V};
    \node[lbl, font=\scriptsize, anchor=east] at (-3, -1.5) 
      {8 slots:};
    \node[font=\scriptsize, text=MedGray] at (1.2, 0) {slot 3};

    \node[font=\scriptsize, text=MedGray] at (-2.5, -1.5) {slot 6};
    \node[font=\scriptsize, text=MedGray] at (2.5, -1.5) {slot 7};
    
    % Level 2: 4 copies
    \node[voidslot, minimum width=0.7cm] (v2a) at (-3, -3.2) {V};
    \node[voidslot, minimum width=0.7cm] (v2b) at (-1, -3.2) {V};
    \node[voidslot, minimum width=0.7cm] (v2c) at (1, -3.2) {V};
    \node[voidslot, minimum width=0.7cm] (v2d) at (3, -3.2) {V};
    \node[lbl, font=\scriptsize, anchor=east] at (-4.2, -3.2) 
      {16 slots:};
    \node[font=\scriptsize, text=MedGray] at (-3, -3.9) {12};
    \node[font=\scriptsize, text=MedGray] at (-1, -3.9) {13};
    \node[font=\scriptsize, text=MedGray] at (1, -3.9) {14};
    \node[font=\scriptsize, text=MedGray] at (3, -3.9) {15};
    
    % Arrows
    \draw[arrow, AccentRed] (v0.south) -- (v1a.north);
    \draw[arrow, AccentRed] (v0.south) -- (v1b.north);
    \draw[arrow, AccentRed] (v1a.south) -- (v2a.north);
    \draw[arrow, AccentRed] (v1a.south) -- (v2b.north);
    \draw[arrow, AccentRed] (v1b.south) -- (v2c.north);
    \draw[arrow, AccentRed] (v1b.south) -- (v2d.north);
    
    % Count annotation
    \node[lbl, font=\small, align=center] at (6, -1.5) 
      {$2^0 = 1$ copy\\$2^1 = 2$ copies\\$2^2 = 4$ copies\\
       $\vdots$\\$2^d$ copies after\\$d$ expansions};
  \end{tikzpicture}
  \caption{Void duplicates grow exponentially: each expansion doubles 
           every existing copy. However, the table also doubles, so 
           the \emph{fraction} of void entries remains constant per 
           generation.}
  \label{fig:void-tree}
\end{figure}

\subsection{Why Duplication Is Correct}

A void entry has zero fingerprint bits. When a query encounters a 
void entry in its target run, the comparison is vacuously true: 
there are no bits to disagree on. Therefore, a void entry matches 
\emph{any} query to its slot. This property ensures correctness in 
two ways:

\begin{enumerate}
  \item \textbf{No false negatives:} The original element's hash 
        determines exactly one of the two candidate slots. Since the 
        void entry occupies \emph{both}, it is guaranteed to be found 
        regardless of which slot the query reaches.
  \item \textbf{False positives are bounded:} A void entry in the 
        ``wrong'' slot may cause additional false positives. However, 
        since the void entry would have matched any query anyway (zero 
        bits to compare), the extra copy does not change the per-query 
        false positive probability. It only affects which queries 
        encounter it.
\end{enumerate}

\subsection{How Queries Find Void Entries}

It is important to understand why duplication guarantees that a query 
always finds the right slot despite the entry being void. At query 
time, the caller has access to the original key and can hash it fresh. 
The hash function produces a full output (e.g., 64 bits), regardless 
of how many bits the filter stored internally. The query computes 
$\log_2(\text{table size})$ bits for the quotient and goes directly 
to that one canonical slot. The query never needs to ``search'' 
across all the duplicates; it knows exactly where to look because it 
has the full hash.

The duplication exists for a different reason: during \emph{expansion}, 
the filter does not have the original key. It only has the stored 
fingerprint. When that fingerprint is empty (void), the filter cannot 
decide which child slot the entry belongs to, so it places the entry 
in both. This ensures that whichever slot the full hash eventually 
resolves to at query time, a void entry will be waiting there.

\subsection{Why Duplication Doesn't Explode}

The exponential growth of void copies (Figure~\ref{fig:void-tree}) 
appears alarming. However, two exponential effects work against each 
other and cancel exactly.

Older generations (entries inserted long ago) need the \emph{most} 
duplicates per entry, because they have been void through many 
expansions. But older generations also contributed the \emph{fewest} 
entries, because the table was much smaller when they were inserted. 
Conversely, recently void generations contributed many entries but 
need very few duplicates (perhaps just one). These two effects, 
the exponential growth in duplicates and the exponential shrinkage 
in generation size, offset each other precisely.

To derive this formally, consider the table starting at size $C_0$. 
Each generation doubles the table, so Generation $j$ sees a table 
of size $C_0 \cdot 2^j$ and Generation $X$ (the current one) sees 
size $C_0 \cdot 2^X$. Generation $j$ added roughly half the table's 
entries at the time of its expansion. The number of new entries 
contributed by Generation $j$ is approximately $C_0 \cdot 2^{j-1}$ 
(the difference between the table size at Generation $j$ and 
Generation $j-1$). The fraction of total entries that came from 
Generation $j$ is therefore:
\[
  f(j) = \frac{C_0 \cdot 2^{j-1}}{C_0 \cdot 2^X} = 2^{-X+j-1}
\]
as shown in Equation 1 from \cite{aleph}.

An entry from Generation $j$ has survived $X - j$ expansions. Each 
expansion consumes one fingerprint bit for the expanded quotient. 
The entry becomes void when all $F$ fingerprint bits are consumed, 
which occurs when $X - j \geq F$. After that point, the entry has 
been void for $(X - j) - F$ additional expansions. During each of 
those additional expansions, the void entry is duplicated into both 
child slots, doubling the number of copies. The total number of 
duplicates per void entry is therefore $2^{X - j - F}$.

The fraction of slots occupied by void duplicates from Generation 
$j$ is the product of how many entries Generation $j$ contributed 
and how many duplicates each entry has:
\begin{equation}
  \alpha \cdot f(j) \cdot 2^{X-j-F} 
  = \alpha \cdot 2^{-X+j-1} \cdot 2^{X-j-F} 
  = \alpha \cdot 2^{-F-1}
  \label{eq:void-fraction}
\end{equation}

The $(X - j)$ terms cancel completely. Every void generation, 
regardless of age, contributes exactly the same fraction of the 
table: $\alpha \cdot 2^{-F-1}$. This balancing act is the key insight. The 
entries requiring the most duplication are exponentially rare, and 
that exponential rarity perfectly offsets the exponential duplication.

The total fraction of void entries is then this constant multiplied 
by the number of void generations. A generation is void when 
$X - j \geq F$, so $j$ ranges from $0$ to $X - F$, giving 
$X - F + 1$ void generations:
\begin{equation}
  \gamma(X) = \alpha \cdot 2^{-F-1} \cdot (X - F + 1)
  \label{eq:total-void-fraction}
\end{equation}

This grows only \emph{linearly} with $X$ (the number of expansions), 
and $2^{-F-1}$ is extremely small for any reasonable fingerprint 
length. A factor of $2^{-F-1}$ with $F = 10$ means each void 
generation contributes only $1/2048$ of the table.

\begin{examplebox}[Void Fraction in Practice]
With $F = 10$ (initial fingerprint length) and $X = 20$ (number of 
expansions), the total void fraction is:
\[
  \gamma(20) = 0.8 \cdot 2^{-11} \cdot (20 - 10 + 1) 
  = 0.8 \cdot \frac{11}{2048} \approx 0.43\%
\]
Less than half a percent of the table is occupied by void duplicates 
after 20 expansions. The table can hold over a million entries at 
this point, so the overhead is negligible.
\end{examplebox}


% ============================================================================
\section{Constant-Time Queries}
\label{sec:constant-queries}
% ============================================================================

With void duplication in place, we can now see why queries are 
$\bigO{1}$. The key consequence is that \textbf{all entries, 
including void entries, reside in a single main hash table}. This 
eliminates the need to search secondary tables during queries.

\textbf{Query algorithm.} To query for element $x$:
\begin{enumerate}
  \item Compute the full hash $h(x)$ using the hash function. This 
        produces a full-length output (e.g., 64 bits), regardless of 
        how many bits are stored internally.
  \item Extract the quotient $q$ (the bottom $\log_2(\text{table 
        size})$ bits) and the remainder $r$ (the next $F$ bits).
  \item Go to canonical slot $q$ in the main hash table.
  \item If the slot is empty or no matching entry exists in the run, 
        return \textsc{no} (definitely not in the set).
  \item If a matching fingerprint is found, return \textsc{yes} 
        (probably in the set). If a void entry is found in the run, 
        also return \textsc{yes}, since a void entry has zero bits to 
        compare against and therefore matches any query trivially.
\end{enumerate}

Every query accesses only the main hash table and terminates in 
$\bigO{1}$ time. Compare this to InfiniFilter, where a query for 
a non-existing or old key must traverse up to 
$\bigO{\log(N)/F}$ hash tables in the Fixed-Width Regime 
\cite{infinifilter}.

\begin{figure}[H]
  \centering
  \begin{tikzpicture}[scale=0.9, every node/.style={transform shape}]
    % InfiniFilter query path
    \node[lbl, font=\bfseries] at (-3, 1) {InfiniFilter query:};
    
    \node[slot, minimum width=1.8cm, fill=BitOne] (m1) at (-3, 0) 
      {\scriptsize main};
    \node[slot, minimum width=1.4cm, fill=orange!10] (s1) at (-3, -1.2) 
      {\scriptsize secondary};
    \node[slot, minimum width=1cm, fill=red!10] (a1) at (-3, -2.4) 
      {\scriptsize aux.};
    
    \draw[arrow, AccentRed] (m1.south) -- (s1.north) 
      node[midway, right, font=\scriptsize] {miss};
    \draw[arrow, AccentRed] (s1.south) -- (a1.north) 
      node[midway, right, font=\scriptsize] {miss};
    \node[lbl, text=AccentRed, font=\scriptsize] at (-3, -3.2) 
      {$\bigO{\log N}$ tables};
    
    % Aleph query path
    \node[lbl, font=\bfseries] at (3, 1) {Aleph Filter query:};
    
    \node[slot, minimum width=1.8cm, fill=BitOne] (m2) at (3, 0) 
      {\scriptsize main (has everything)};
    
    \draw[arrow, AccentGreen, very thick] (3, 0.8) -- (m2.north);
    \node[lbl, text=AccentGreen, font=\scriptsize] at (3, -1) 
      {$\bigO{1}$ always};
    
    \node[slot, minimum width=1.4cm, fill=gray!5] (s2) at (3, -2) 
      {\scriptsize secondary};
    \node[font=\scriptsize, text=MedGray, align=center] at (3, -2.8) 
      {only for deletes,\\never for queries};
  \end{tikzpicture}
  \caption{InfiniFilter queries may traverse multiple tables; Aleph 
           Filter queries always access only the main table.}
  \label{fig:query-comparison}
\end{figure}


% ============================================================================
\section{Tombstone-Based Deletes}
\label{sec:tombstone-deletes}
% ============================================================================

Void duplication solves the query problem, but it introduces a 
complication for deletes. When a void entry is deleted, \emph{all of 
its duplicates} must also be removed to prevent the filter from 
returning false positives for the deleted key. A void entry may 
have $2^{k-b}$ duplicates (where $k = \log_2(\text{table size})$ and 
$b$ is the length of its mother hash from 
Section~\ref{sec:mother-hash}), and removing them all immediately 
could be expensive \cite{aleph}.

The Aleph Filter addresses this with a \textbf{two-phase lazy 
deletion} scheme.

\subsection{Phase 1: Immediate Tombstone ($\bigO{1}$)}

When deleting a key whose only matching entry in the target run is a 
void entry:

\begin{enumerate}
  \item Replace the void entry at the canonical slot with a 
        \textbf{tombstone}: a special bit pattern of all 1s 
        (Figure~\ref{fig:slot-format}).
  \item Add the canonical slot address to a \textbf{deletion queue} 
        (an append-only array).
\end{enumerate}

The tombstone ensures that subsequent queries for the deleted key 
return \textsc{no}, providing immediate correctness. Since the void 
entry previously matched everything at that slot (zero fingerprint 
bits means vacuously matching any query), the tombstone replaces that 
blanket ``yes'' with a definitive ``no.'' This phase takes $\bigO{1}$ 
time.

\subsection{Using the Mother Hash for Duplicate Removal}

Recall from Section~\ref{sec:mother-hash} that the mother hash is the 
accumulated slot address of a void entry, containing all the hash 
bits ever stored for it. The mother hash is stored in a 
\textbf{secondary hash table}, which is itself a smaller quotient 
filter. Because the secondary table has fewer slots, it requires 
fewer quotient bits, so the mother hash can be split into a short 
quotient and a long fingerprint. This means the entry that was void 
in the main table (zero fingerprint bits) has a meaningful fingerprint 
again in the secondary table.

\subsection{Phase 2: Deferred Duplicate Removal}

Right before the next expansion, the Aleph Filter processes the 
deletion queue:

\begin{enumerate}
  \item Pop a canonical slot address from the queue.
  \item Look up the address in the \textbf{secondary hash table} to 
        find the void entry's mother hash.
  \item From the mother hash, compute the locations of all void 
        duplicates. If the mother hash is $b$ bits and the table has 
        $2^k$ slots, then the entry has $2^{k-b}$ duplicates. Their 
        slot addresses share the mother hash as the bottom $b$ bits, 
        with all $2^{k-b}$ permutations of the remaining upper bits:
        \[
          \text{duplicate slots} = \{s \in \{0, \dots, 2^k - 1\} 
          : s \mod 2^b = \text{mother hash}\}
        \]
        The mother hash forms the lower bits of each duplicate's slot 
        address. The upper bits cycle through every possibility, since 
        those are precisely the bits the filter had no information 
        about (the bits beyond the void threshold).
  \item Remove a void entry from each of these canonical slots.
  \item Remove the mother hash from the secondary hash table.
\end{enumerate}

\begin{examplebox}[Deletion Walkthrough (from \cite{aleph}, Figure 9)]
Delete Key $X$ with $h(X) = 001101$. The main table has $2^3 = 8$ 
slots.

\textbf{Phase 1:}
\begin{enumerate}
  \item Find Key $X$'s canonical slot: 101 (from the bottom 3 bits of 
        the hash, i.e.\ the quotient).
  \item The only matching entry is a void entry. Replace it with a 
        tombstone.
  \item Add ``101'' to the deletion queue.
\end{enumerate}

\textbf{Phase 2} (before next expansion):
\begin{enumerate}
  \item Pop ``101'' from the queue.
  \item Search the secondary hash table for the longest matching 
        mother hash. Find entry at Slot 1 with fingerprint 0. 
        Concatenate: longest matching mother hash $= 01$ ($b = 2$ bits).
  \item Table has $2^3 = 8$ slots ($k = 3$). Number of duplicates: 
        $2^{3-2} = 2$.
  \item Duplicate locations: bottom 2 bits $= 01$, permute the top 
        bit: slots $\texttt{001}$ and $\texttt{101}$.
  \item Remove a void entry from each slot. Remove the mother hash 
        from the secondary table.
\end{enumerate}
\end{examplebox}

\subsection{Why Delete the Longest Matching Mother Hash?}

When multiple void entries share the same canonical slot, there may 
be multiple matching mother hashes of different lengths in the 
secondary table. The Aleph Filter always removes the entry with the 
\textbf{longest matching mother hash} (i.e., the fewest duplicates) 
\cite{aleph}.

To see why, consider two void entries $X$ and $Y$ that share a 
canonical slot. $X$ is older, so its mother hash is shorter (say 2 
bits) and it has 8 duplicates. $Y$ is newer, so its mother hash is 
longer (say 4 bits) and it has only 2 duplicates. $Y$'s duplicate 
slots are a \emph{subset} of $X$'s duplicate slots, because $Y$'s 
mother hash is a more specific prefix that narrows down to fewer 
locations.

If we were to delete the shorter mother hash ($X$'s), we would 
remove void entries from all 8 of $X$'s slots. But $Y$ only has 
duplicates at 2 of those 8 slots. The other 6 slots that we just 
cleared contained void entries that $Y$'s queries might need. A 
future query for $Y$ could land on one of the 6 slots that no longer 
has a void entry and return a false negative, which filters must 
never produce.

By removing the longest matching mother hash ($Y$'s, with fewer 
duplicates), we only remove 2 void entries. $X$'s 8 duplicates 
remain intact. Since $Y$'s slots were a subset of $X$'s, every 
slot that $Y$ occupied is still covered by $X$. No false negatives 
are possible.

The rule is: remove the entry with the most specific (longest) mother 
hash first, because it disturbs the fewest slots and the remaining 
entry's broader set of duplicates continues to cover everything.

\subsection{Computational Cost}

The tombstone placement and queue insertion in Phase 1 take $\bigO{1}$.
The duplicate removal in Phase 2 is deferred to the next expansion. 
The cost of Phase 2 is spread across all the insertions that triggered 
the expansion. Formally, the total number of void duplicates in the 
Fixed-Width Regime is at most $\bigO{2^{-F} \cdot X \cdot N}$ 
(Section 4.2 of \cite{aleph}). Since this work is performed after $N$ 
insertions, the amortized cost per insertion is sub-constant as long 
as $X < 2^F$, which is the maximum number of supported expansions. 

To understand what ``amortized $\bigO{1}$'' means here: most 
individual deletions are cheap (just place a tombstone). Occasionally, 
the deferred Phase 2 cleanup during an expansion is more expensive. 
But when that rare expensive cost is averaged across all the cheap 
operations leading up to it, the average per-operation cost remains 
$\bigO{1}$.


% ============================================================================
\section{Rejuvenation}
\label{sec:rejuvenation}
% ============================================================================

The tombstone mechanism handles deletes. The third and final 
mechanism, \textbf{rejuvenation}, handles a complementary problem: 
keeping the FPR low over time by refreshing old entries.

When a query returns a \emph{true positive} (the queried key actually 
exists in the underlying data set), the application typically fetches 
the full key from storage. At this point, the Aleph Filter can 
\textbf{rejuvenate} the entry: rehash the full key to obtain a fresh, 
full-length fingerprint and replace the short or void entry with this 
new fingerprint \cite{aleph}.

Rejuvenation reduces the FPR by replacing entries that contribute 
high false positive probabilities (short or void fingerprints) with 
entries that contribute low probabilities (full-length fingerprints). 
In essence, the entry is ``reborn'' with a fresh $F$-bit fingerprint 
as if it had just been inserted, undoing all the damage caused by 
successive expansions.

\subsection{Rejuvenation Procedure}

When a query finds that the only matching entry in the target run is 
a void entry, and the query is a true positive:

\begin{enumerate}
  \item \textbf{Immediate:} Replace the void entry at the canonical 
        slot with the fresh full-length fingerprint computed from the 
        fetched key. Add the canonical slot address to a 
        \textbf{rejuvenation queue}.
  \item \textbf{Deferred:} Before the next expansion, pop addresses 
        from the rejuvenation queue. For each, look up the longest 
        matching mother hash in the secondary hash table, compute the 
        locations of all void duplicates, and remove them. This is 
        identical to the delete procedure 
        (Section~\ref{sec:tombstone-deletes}), except the void entry 
        at the queried key's canonical slot has already been replaced 
        with a real fingerprint rather than a tombstone.
\end{enumerate}

\begin{examplebox}[Rejuvenation Walkthrough (from \cite{aleph}, Figure 11)]
Rejuvenate Key $X$ with $h(X) = 001100$. Main table has $2^3 = 8$ 
slots.

\begin{enumerate}
  \item Query finds void entry at canonical slot 100.
  \item Application fetches Key $X$ from storage (true positive).
  \item Rehash Key $X$ to get a fresh fingerprint. Replace the void 
        entry at slot 100 with fingerprint \texttt{0001}.
  \item Add ``100'' to the rejuvenation queue.
  \item Before next expansion: look up ``100'' in the secondary table. 
        Find longest matching mother hash $= 00$ ($b = 2$ bits). 
        Duplicates are at slots $\texttt{000}$ and $\texttt{100}$.
  \item Since slot 100 was already rejuvenated (it now has a real 
        fingerprint), only the duplicate at slot $\texttt{000}$ needs 
        to be removed.
  \item Remove the mother hash from the secondary table.
\end{enumerate}
\end{examplebox}

Rejuvenation is only effective when queries target older entries 
(those with short or void fingerprints). In workloads that 
predominantly query recent entries, rejuvenation has little 
opportunity to help. The Widening and Predictive Regimes 
(Section~\ref{sec:regimes}) provide complementary mechanisms that 
scale the FPR without relying on the query workload.


% ============================================================================
\section{FPR Analysis}
\label{sec:fpr-analysis}
% ============================================================================

With all three mechanisms in place (void duplication, tombstones, and 
rejuvenation), we can now analyze the false positive rate. The key 
result is that \textbf{every generation contributes equally to the 
FPR}, whether its entries are void or not \cite{aleph}.

\subsection{Per-Generation FPR}

Let $f(j) \approx 2^{-X+j-1}$ denote the fraction of entries from 
Generation $j$, where $X$ is the current generation.

\textbf{Non-void entries} (Generation $j$ where $X - j < F$): These 
entries have fingerprints of $F - (X - j)$ bits remaining. Each 
expansion steals one fingerprint bit, so the collision probability 
doubles (the fingerprint halves in distinguishing power). But the 
generation's share of the table also halves with each expansion. 
These two effects cancel:
\[
  \alpha \cdot f(j) \cdot 2^{-F+(X-j)} = \alpha \cdot 2^{-F-1}
\]

\textbf{Void entries} (Generation $j$ where $X - j \geq F$): These 
entries match any query with probability 1 (zero fingerprint bits 
means no basis for rejection). However, their duplication keeps 
their fraction constant. The generation's small size and large 
duplicate count cancel as shown in 
Equation~\ref{eq:void-fraction}:
\[
  \alpha \cdot f(j) \cdot 2^{X-j-F} \cdot 1 = \alpha \cdot 2^{-F-1}
\]

The same value in both cases. For non-void entries, fingerprint 
shrinkage increases collision probability but the generation is 
proportionally smaller. For void entries, the match probability is 1 
but the duplication count and generation size cancel to the same 
constant. Every generation, void or not, contributes exactly 
$\alpha \cdot 2^{-F-1}$ to the FPR.

\subsection{Total FPR}

Summing across all $X + 1$ generations (indexed $0$ through $X$) 
plus the current generation:

\begin{equation}
  \text{FPR} \approx \sum_{j=0}^{X+1} \alpha \cdot 2^{-F-1} 
  \lesssim \alpha \cdot (\log_2(N) + 2) \cdot 2^{-F-1}
  \label{eq:aleph-fpr-fixed}
\end{equation}

This matches InfiniFilter's FPR in the Fixed-Width Regime 
(Equation 2 from \cite{aleph}). The Aleph Filter does not increase 
the FPR by duplicating void entries. It merely changes \emph{where} 
they are stored, not how many effective false positives they produce.

\subsection{Widening Regime}

In the Widening Regime, entries in Generation $j$ are assigned 
fingerprints of $\ell(j) = F + 2 \cdot \log_2(j + 1)$ bits. The 
extra bits compensate for the longer lifetime these entries will 
have, since they will survive more expansions and lose more bits. 
The FPR converges using the Basel series identity 
$\sum_{j=1}^{\infty} j^{-2} = \pi^2/6$:

\begin{equation}
  \text{FPR} \approx \alpha \cdot 2^{-F-1} \cdot 
  \sum_{j=0}^{\infty} \frac{1}{(j+1)^2} 
  \lesssim \alpha \cdot 2^{-F-1} \cdot \frac{\pi^2}{6} 
  \lesssim \alpha \cdot 2^{-F}
  \label{eq:aleph-fpr-widening}
\end{equation}

The FPR converges to a constant $\bigO{2^{-F}}$, independent of how 
much the data grows. This is achieved at the cost of $F + \bigO{\log 
\log N}$ bits per entry \cite{aleph}.


% ============================================================================
\section{Operating Regimes}
\label{sec:regimes}
% ============================================================================

The Aleph Filter supports three operating regimes, each offering a 
different tradeoff between FPR, memory, and maximum expansions:

\begin{table}[H]
  \centering
  \caption{Aleph Filter operating regimes \cite{aleph}. $N$ is the data 
           size relative to initial capacity, $F$ is the initial 
           fingerprint length, and $N_{est}$ is the estimated final 
           data size.}
  \label{tab:regimes}
  \begin{tabular}{l c c c}
    \hline
    \textbf{Regime} & \textbf{FPR} & \textbf{Bits/entry} & 
    \textbf{Max expansions} \\
    \hline
    Fixed-Width & $\bigO{2^{-F} \cdot \log N}$ & $F$ & $2^F$ \\
    Widening & $\bigO{2^{-F}}$ & $F + \bigO{\log \log N}$ & $\infty$ \\
    Predictive & $\bigO{2^{-F}}$ & $F + \bigO{\log|\log(N/N_{est})|}$ 
    & $\infty$ \\
    \hline
  \end{tabular}
\end{table}

The \textbf{Predictive Regime} is a novel contribution of the Aleph 
Filter paper \cite{aleph}. Given an estimate $N_{est}$ of how much the 
data will grow, it assigns \emph{longer} fingerprints initially and 
\emph{shorter} ones as the data approaches the estimate. When the 
data size reaches $N_{est}$, all fingerprints have shrunk to exactly 
$F$ bits, achieving an FPR vs.\ memory tradeoff on par with static 
filters.

If the data exceeds the estimate, the Predictive Regime transitions 
to Widening Regime behavior, assigning increasingly longer 
fingerprints. The memory cost is $F + \bigO{\log|\log(N/N_{est})|}$, 
which is significantly better than the Widening Regime's 
$F + \bigO{\log \log N}$ when $N \approx N_{est}$.


% ============================================================================
\section{Architecture Summary}
% ============================================================================

The complete Aleph Filter architecture consists of:

% ---- Diagram: Architecture ----
\begin{figure}[H]
  \centering
  \begin{tikzpicture}
    % Main table
    \node[draw, rounded corners, fill=BitOne, minimum width=7cm, 
          minimum height=1.5cm] (main) at (0, 0) {};
    \node[font=\bfseries] at (0, 0.3) {Main Hash Table};
    \node[font=\scriptsize, align=center] at (0, -0.2) 
      {Normal entries + void duplicates + tombstones\\
       \textbf{All queries go here only}};
    
    % Secondary table
    \node[draw, rounded corners, fill=orange!10, minimum width=4.8cm, 
          minimum height=2.0cm] (sec) at (0, -2.3) {};
    \node[font=\bfseries, font=\small] at (0, -2.0) {Secondary Hash Table};
    \node[font=\scriptsize, align=center] at (0, -2.6) 
      {Mother hashes of void entries\\
       For deletes and rejuvenation only};
    
    % Auxiliary tables
    \node[draw, rounded corners, fill=red!8, minimum width=4cm, 
          minimum height=1.3cm] (aux) at (0, -4.4) {};
    \node[font=\bfseries, font=\small] at (0, -4.4) {Auxiliary Tables};
    \node[font=\scriptsize] at (0, -4.70) {Chain of older sealed tables};
    
    % Queues
    \node[draw, rounded corners, fill=AccentTeal!10, minimum width=2.8cm, 
          minimum height=0.7cm] (dq) at (5.5, -1.2) {};
    \node[font=\scriptsize] at (5.5, -1.2) {Deletion Queue};
    
    \node[draw, rounded corners, fill=AccentGreen!10, minimum width=2.8cm, 
          minimum height=0.7cm] (rq) at (5.5, -3.2) {};
    \node[font=\scriptsize] at (5.5, -3.2) {Rejuvenation Queue};
    
    % Arrows
    \draw[arrow] (main.south) -- (sec.north) 
      node[midway, right=4pt, font=\scriptsize] 
      {void mother hashes};
    \draw[arrow] (sec.south) -- (aux.north) 
      node[midway, right=4pt, font=\scriptsize] {when full};
    \draw[arrow] (main.east) -- (dq.west) 
      node[midway, above, right=2pt, yshift=7pt, font=\scriptsize] {delete};
    \draw[arrow] (main.east) -- (rq.west) 
      node[midway, below, xshift=30pt, yshift=-20pt, font=\scriptsize] {rejuvenate};
    \draw[dasharrow, AccentRed] (dq.south) -- (sec.east) 
      node[midway, right=10pt, yshift=-4pt, font=\scriptsize, align=left] 
      {before next expansion};
    \draw[dasharrow, AccentRed] (rq.west) -- (sec.east);
    
    % Query label
    \node[hashbox] (q) at (-5.5, 0) {query $x$};
    \draw[arrow, AccentGreen, very thick] (q.east) -- (main.west)
      node[midway, above, font=\scriptsize, text=AccentGreen] 
      {$\bigO{1}$};
  \end{tikzpicture}
  \caption{Aleph Filter architecture. Queries access only the main 
           table. The secondary and auxiliary tables store mother hashes 
           for delete and rejuvenation operations. Queues defer 
           duplicate removal to the next expansion.}
  \label{fig:architecture}
\end{figure}


% ============================================================================
\section{Comparison}
% ============================================================================

\begin{table}[H]
  \centering
  \caption{Complete comparison of filter expansion techniques. The 
           Aleph Filter is the first to achieve $\bigO{1}$ for all 
           operations while supporting infinite expansion with a 
           stable FPR \cite{aleph}.}
  \label{tab:full-comparison}
  \begin{tabular}{l c c c c c}
    \hline
    & \textbf{Query} & \textbf{Delete} & \textbf{Expand} 
    & \textbf{Stable FPR} & \textbf{Insert} \\
    \hline
    Bloom \cite{bloom}
    & $\bigO{k}$ & \texttimes & \texttimes & \texttimes & $\bigO{k}$ \\
    Cuckoo \cite{cuckoo}
    & $\bigO{1}$ & \checkmark & \texttimes & \texttimes & $\bigO{1}$* \\
    Quotient \cite{quotient}
    & $\bigO{1}$* & \checkmark & \checkmark & \texttimes & $\bigO{1}$* \\
    InfiniFilter \cite{infinifilter}
    & $\bigO{\log(N)/F}$ & \checkmark & \checkmark & \checkmark & $\bigO{1}$ \\
    \textbf{Aleph Filter} \cite{aleph}
    & $\bigO{1}$ & \checkmark & \checkmark & \checkmark & $\bigO{1}$ \\
    \hline
  \end{tabular}
  
  \smallskip
  {\small * = amortized, meaning individual operations are usually 
  $\bigO{1}$ but rare expensive operations (eviction chains or 
  run shifts) bring the average to $\bigO{1}$ when spread across 
  many operations. Performance degrades at high load factors. For
  InfiniFilter, $\bigO{\log(N)/F}$ is equivalent to $\bigO{\log N}$ when
  $F$ is treated as a constant.}
\end{table}

Each row in Table~\ref{tab:full-comparison} solves a limitation of 
the row above it: Bloom filters cannot delete, cuckoo filters cannot 
expand, quotient filters cannot maintain a stable FPR, and 
InfiniFilter's queries degrade as data grows. The Aleph Filter 
achieves all properties simultaneously: constant-time operations 
that remain stable no matter how much the data grows.
\chapter{Rust Implementation}

This chapter walks through our Rust implementation of the Aleph Filter. We begin with the core data structures, then explain each operation and how it maps to the algorithms described in the previous chapters. Finally, we compare our implementation against the official Java reference to clarify what matches and what was intentionally simplified.

\section{Data Structures}

The filter is built on three core types: a \texttt{Slot} for storing fingerprints, a \texttt{SlotMetadata} for tracking the three quotient filter flags, and the \texttt{AlephFilter} struct that ties everything together.

\begin{lstlisting}[caption={The main AlephFilter struct}]
pub struct AlephFilter {
    slots: Vec<Slot>,             // fingerprint storage
    metadata: Vec<SlotMetadata>,  // per-slot O/C/S flags
    num_slots: usize,             // logical slots (power of 2)
    num_extension_slots: usize,   // overflow buffer
    quotient_bits: u32,           // log2(num_slots)
    base_fp_bits: u32,            // fingerprint width at creation
    num_expansions: usize,        // times we have doubled
    num_items: usize,             // items currently stored
    max_load_factor: f64,         // expansion threshold (0.8)
}
\end{lstlisting}

Each \texttt{Slot} packs a fingerprint value and its bit-length into a single 64-bit integer: the lower 56 bits hold the fingerprint, and the upper 8 bits hold the length. Two special sentinel values are reserved: a \texttt{Void} marker for exhausted fingerprints, and a \texttt{Tombstone} marker for deleted void entries.

Each \texttt{SlotMetadata} is a single byte with three bit-flags:
\begin{itemize}
    \item \textbf{Occupied} (bit 0): this canonical slot has a run associated with it.
    \item \textbf{Continuation} (bit 1): this slot continues a run that started earlier.
    \item \textbf{Shifted} (bit 2): this slot's data was displaced from its canonical position.
\end{itemize}

These three flags are exactly the standard quotient filter metadata described in Chapter~2. Together, they allow us to locate any element's run in $O(1)$ expected time by walking backward to find the cluster start and then forward to find the target run.

% ============================================================================
\section{Operations}
% ============================================================================

With the data structures in place, we now describe each operation and how it maps to the algorithms from the previous chapters.

% Subsection
\subsection{Insertion}

Insertion in our Rust code is split into a public API step (hash split + slot construction) and an internal quotient-filter step (new run vs existing run).

\begin{lstlisting}[caption={Insertion path in our Rust implementation}]
pub fn insert(&mut self, key: &[u8]) {
  if self.load_factor() >= self.max_load_factor {
    self.expand();
  }
  let hash = hash_key(key);
  let r = self.fp_bits();
  let (slot_idx, fp) = split_hash(hash, self.quotient_bits, r);
  let canonical = slot_idx as usize;
  let slot = if r == 0 { Slot::void_marker() } else { Slot::new(fp, r as u8) };
  self.qf_insert_slot(canonical, slot);
  self.num_items += 1;
}

fn qf_insert_slot(&mut self, canonical: usize, slot: Slot) -> bool {
  let does_run_exist = self.metadata[canonical].is_occupied();
  if !does_run_exist {
    return self.insert_new_run(canonical, slot);
  } else {
    let run_start = self.find_run_start(canonical);
    return self.insert_into_run(slot, run_start);
  }
}
\end{lstlisting}

This matches the quotient-filter model: split into quotient/remainder, then either create a new run or append into an existing run while preserving \texttt{occupied}, \texttt{continuation}, and \texttt{shifted} metadata.

% Subsection
\subsection{Query}

Query logic in the implementation is also two-stage: compute hash parts in \texttt{contains}, then run-local matching in \texttt{qf\_search}/\texttt{find\_in\_run}.

\begin{lstlisting}[caption={Query path in our Rust implementation}]
pub fn contains(&self, key: &[u8]) -> bool {
  let hash = hash_key(key);
  let r = self.fp_bits();
  let (slot_idx, fp) = split_hash(hash, self.quotient_bits, r);
  let canonical = slot_idx as usize;
  return self.qf_search(fp, canonical);
}

fn qf_search(&self, fp: u64, canonical: usize) -> bool {
  if !self.metadata[canonical].is_occupied() {
    return false;
  }
  let run_start = self.find_run_start(canonical);
  return self.find_in_run(run_start, fp);
}

fn find_in_run(&self, start: usize, fp: u64) -> bool {
  let mut pos = start;
  let r = self.fp_bits();
  loop {
    if self.slots[pos].is_tombstone() {
      // skip tombstones
    } else if self.slots[pos].matches(fp, r as u8) {
      return true;
    }
    pos += 1;
    if pos >= self.total_slots() || !self.metadata[pos].is_continuation() {
      break;
    }
  }
  return false;
}
\end{lstlisting}

Behavioral detail: tombstones are skipped, void entries can still match, and comparison uses the current effective fingerprint width.

% Subsection
\subsection{Expansion}

The Rust implementation expands by iterating existing entries, doubling the table, then reinserting with pivot-bit routing. Void entries are duplicated to both candidate buckets.

\begin{lstlisting}[caption={Expansion core in our Rust implementation}]
fn expand(&mut self) {
  let entries = self.iterate_entries();
  let old_q_bits = self.quotient_bits;

  self.num_slots *= 2;
  self.quotient_bits += 1;
  self.num_expansions += 1;
  self.num_extension_slots += 2;
  self.slots = vec![Slot::empty(); self.total_slots()];
  self.metadata = vec![SlotMetadata::new(); self.total_slots()];
  self.num_items = 0;

  let mut void_slots: std::collections::HashSet<usize> = std::collections::HashSet::new();

  for (bucket, fingerprint, fp_len, is_void, _old_pos) in entries {
    if is_void {
      let s0 = bucket;
      let s1 = bucket | (1 << old_q_bits);
      if s0 < self.num_slots && void_slots.insert(s0) {
        self.qf_insert_slot(s0, Slot::void_marker());
        self.num_items += 1;
      }
      if s1 < self.num_slots && void_slots.insert(s1) {
        self.qf_insert_slot(s1, Slot::void_marker());
        self.num_items += 1;
      }
    } else if fp_len > 0 {
      let pivot = fingerprint & 1;
      let new_bucket = bucket | ((pivot as usize) << old_q_bits);
      let new_fp = fingerprint >> 1;
      let new_len = fp_len - 1;
      let slot = if new_len == 0 { Slot::void_marker() } else { Slot::new(new_fp, new_len) };
      self.qf_insert_slot(new_bucket, slot);
      self.num_items += 1;
    }
  }
}
\end{lstlisting}

This is the key Aleph behavior in code: one bit is sacrificed per expansion for normal entries, while void entries are duplicated to preserve no-false-negatives lookup behavior.

\subsubsection{Void Entries}

After enough expansions, some entries exhaust all their fingerprint bits ($r = 0$). These become \textit{void entries} as they have no bits left to pivot on, so we cannot determine which half of the expanded table they belong to. The Aleph Filter's solution is to \textbf{duplicate} the void entry into \textit{both} possible new slots:

\begin{align}
    s_0 &= b \\
    s_1 &= b \lor (1 \ll q)
\end{align}

This duplication ensures that queries to either new bucket will find the void marker and return \texttt{true}, preserving the no-false-negatives guarantee. The cost is a small amount of space redundancy, but it avoids the $O(\log N)$ secondary table lookups that InfiniFilter requires.

% Subsection
\subsection{Deletion}

Deletion uses full cluster-aware compaction for normal entries, and a tombstone fast path for void entries.

\begin{lstlisting}[caption={Deletion path in our Rust implementation}]
pub fn delete(&mut self, key: &[u8]) -> bool {
  let hash = hash_key(key);
  let r = self.fp_bits();
  let (slot_idx, fp) = split_hash(hash, self.quotient_bits, r);
  let canonical = slot_idx as usize;
  if self.qf_delete(fp, canonical) {
    self.num_items = self.num_items.saturating_sub(1);
    return true;
  } else {
    return false;
  }
}

fn qf_delete(&mut self, fp: u64, canonical: usize) -> bool {
  // ... locate match in run ...
  if self.slots[del_pos].is_void() {
    self.slots[del_pos] = Slot::tombstone();
    return true;
  }
  // ... full cluster compaction for regular entries ...
}
\end{lstlisting}

So the implementation preserves quotient-filter cluster invariants for regular deletions, but intentionally uses tombstone-only behavior for void deletions.

\subsubsection{Void Deletion: Our Simplification}
\label{sec:void-delete-tradeoff}

The paper's full deletion algorithm for void entries involves two phases. In Phase~1, the void entry at the queried slot is immediately replaced with a tombstone. In Phase~2, deferred to the next expansion, the filter looks up the void entry's mother hash in a secondary filter chain, computes the locations of all $2^{k-b}$ duplicate copies (where $k = \log_2(\text{table size})$ and $b$ is the length of the mother hash), and removes every one of them.

Our implementation performs \textbf{only Phase~1}. The relevant code is straightforward:

\begin{lstlisting}[caption={Void deletion in our Rust implementation}]
if self.slots[del_pos].is_void() {
    self.slots[del_pos] = Slot::tombstone();
    return true;
}
\end{lstlisting}

We tombstone the void entry at the slot where the query found it, but we do not track down or remove any of its duplicates. This has two concrete consequences:

\begin{table}[H]
  \centering
  \caption{Behavioral differences when deleting a void entry}
  \label{tab:void-delete}
  \begin{tabular}{p{4.5cm} p{4cm} p{4cm}}
    \hline
    \textbf{Scenario} & \textbf{Paper (full)} & \textbf{Ours (tombstone only)} \\
    \hline
    Query the deleted key at the tombstoned slot & Returns \textsc{false} \checkmark & Returns \textsc{false} \checkmark \\
    \hline
    Query the deleted key at a \textit{different} duplicate slot & Duplicate was removed; returns \textsc{false} \checkmark & Duplicate still present; returns \textsc{true} (false positive) \\
    \hline
    Space reclamation & All $2^{k-b}$ duplicates freed & Duplicates remain, wasting space \\
    \hline
  \end{tabular}
\end{table}

In other words, tombstoning alone blocks the query at the specific slot where the tombstone was placed, but queries that hash to one of the remaining duplicate slots will still encounter a void entry and return \textsc{true}. This is a \textit{false positive for a deleted key}, which the paper's full algorithm avoids.

\paragraph{Why we accepted this tradeoff.}
The full duplicate-removal mechanism requires three additional components that our implementation does not have:

\begin{enumerate}
    \item A \textbf{secondary filter} (itself a quotient filter) to store the mother hash of each void entry when it is first created during expansion.
    \item A \textbf{\texttt{delete\_duplicates}} function that, given the mother hash length $b$, computes all $2^{k-b}$ duplicate slot addresses and removes a void entry from each one.
    \item A \textbf{deletion queue} that batches tombstoned addresses and processes them in bulk before the next expansion.
\end{enumerate}

Implementing these components would roughly double the codebase (approximately 250 additional lines) and introduce a recursive data structure (the secondary filter can itself overflow, requiring a tertiary filter, and so on). We chose to omit this subsystem for two reasons:

\begin{itemize}
    \item \textbf{Scope}: The secondary chain is an auxiliary subsystem for delete correctness, not part of the core Aleph contribution (void duplication for $O(1)$ queries). Omitting it keeps the implementation focused on demonstrating the main algorithm.
    \item \textbf{Workload}: The tradeoff only manifests when void entries are deleted. An entry becomes void only after it has survived at least $F$ expansions (where $F$ is the initial fingerprint length, typically 7--10). In most practical workloads, entries that old represent a tiny fraction of the data, and deleting them is rare.
\end{itemize}

For workloads that do not delete old entries, this gap never manifests. For workloads that do, the consequence is a bounded increase in false positives for deleted keys, not a correctness violation for other keys.

% ============================================================================
\section{Fidelity to the Paper}
% ============================================================================

This section documents exactly what matches and what was intentionally simplified compared the official rust implementation at \texttt{github.com/nivdayan/AlephFilter} (branch \texttt{aleph\_filter}), to make the scope of this implementation clear. 

\subsection{What Matches Exactly}

Table~\ref{tab:matching} lists every component where our Rust code follows the same algorithm as the Java, verified by side-by-side code comparison.

\begin{table}[H]
  \centering
  \caption{Components matching the Java reference implementation}
  \label{tab:matching}
  \begin{tabular}{l l}
    \hline
    \textbf{Component} & \textbf{Status} \\
    \hline
    Hash splitting (quotient / remainder)        & Exact match \\
    Metadata flags (Occupied, Continuation, Shifted) & Exact match \\
    \texttt{find\_run\_start} algorithm          & Exact match \\
    \texttt{find\_new\_run\_location} algorithm  & Exact match \\
    \texttt{insert\_new\_run} with swap chain    & Exact match \\
    \texttt{insert\_into\_run} (push-all-else)   & Exact match \\
    \texttt{search} / \texttt{find\_in\_run}     & Exact match \\
    Iterator state machine (5 metadata cases)    & Exact match \\
    Pivot-bit expansion (\texttt{expand})         & Exact match \\
    Void duplication to both new canonical slots  & Exact match \\
    Full cluster-aware deletion                   & Exact match \\
    Load factor threshold ($0.8$)                 & Exact match \\
    Extension slots formula ($q \times 2$)        & Exact match \\
    \hline
  \end{tabular}
\end{table}

In short: the quotient filter core (insertion, lookup, deletion), the expansion algorithm (pivot-bit sacrifice with void duplication), and the iterator used during expansion all follow the same logic as the Java implementation.

\subsection{What Does Not Match}

Six aspects of the Java implementation are not reproduced in our Rust code. Table~\ref{tab:differences} summarizes each difference alongside the Java's approach and our alternative.

\begin{table}[H]
  \centering
  \caption{Differences between our Rust implementation and the Java reference}
  \label{tab:differences}
  \begin{tabular}{p{3.2cm} p{4.2cm} p{4.8cm}}
    \hline
    \textbf{Feature} & \textbf{Java Reference} & \textbf{Our Rust} \\
    \hline
    Storage model & Packed bit array; each slot is $(3{+}r)$ bits & 64-bit struct per slot \\
    \hline
    Unary encoding & High fingerprint bits encode entry age via unary prefix & Slot stores explicit length; overlap comparison instead \\
    \hline
    Secondary filter chain & Tracks void entry origins for duplicate cleanup on delete & Not implemented; voids are tombstoned \\
    \hline
    Lazy delete batching & Defers duplicate removal; batch-resolves later & Not implemented \\
    \hline
    Rejuvenation & Replaces a void entry with a fresh fingerprint & Not implemented \\
    \hline
    FP growth strategy & Multiple strategies (\textsc{uniform}, \textsc{polynomial}) & Fixed: shrinks by 1 bit per expansion \\
    \hline
  \end{tabular}
\end{table}

Table~\ref{tab:differences} is the authoritative summary of scope differences. We intentionally keep this section brief to avoid repeating the operation-level explanations above.

In short:
\begin{itemize}
  \item We prioritize clarity over bit-packing by using explicit \texttt{Slot} and \texttt{SlotMetadata} structures.
  \item We use explicit fingerprint lengths instead of unary age encoding.
  \item We implement tombstone-only void deletion (no secondary chain, no lazy duplicate cleanup, no rejuvenation).
  \item We use a fixed shrink strategy (one fingerprint bit per expansion).
\end{itemize}

\begin{insightbox}[Key Takeaway]
Our implementation captures the \textbf{core algorithmic contribution} of the Aleph Filter paper: pivot-bit expansion with void duplication, enabling $O(1)$ queries regardless of how many times the filter has expanded. This is the feature that distinguishes the Aleph Filter from InfiniFilter, and it is fully implemented. The remaining differences are either memory optimizations (packed bitmap), performance optimizations (lazy deletes), or auxiliary subsystems for advanced delete semantics (secondary chain, rejuvenation).
\end{insightbox}

\chapter{Comparison and Conclusion}

  \section{Benchmark Methodology}

  To evaluate our Rust Aleph Filter implementation against established probabilistic
  filter libraries, we constructed a comparative benchmark suite measuring five dimensions:
  insertion throughput, positive lookup (all keys present), negative lookup (no keys present),
  deletion throughput, and false positive rate accuracy. All benchmarks target a 1\% false
  positive rate.

  \subsection{Experimental Setup}

  \begin{itemize}
    \item \textbf{System:} Apple M2 (ARM64), macOS 15.3, Darwin 25.3.0
    \item \textbf{Compiler:} Rust 1.93.0, compiled with \texttt{--release} (optimized)
    \item \textbf{Benchmark framework:} Criterion 0.5 (100 samples per point, statistical analysis)
    \item \textbf{Secondary validation:} Custom runner with median-of-5 timing
    \item \textbf{Workload sizes:} $N \in \{1{,}000,\; 10{,}000,\; 100{,}000\}$
    \item \textbf{Keys:} Sequential byte strings (\texttt{"key\_0"}, \texttt{"key\_1"}, \ldots)
  \end{itemize}

  \subsection{Competing Implementations}

  We benchmarked the following Rust crate libraries:

  \begin{itemize}
    \item \textbf{Bloom Filter}: \texttt{bloomfilter} v1.0.16 (standard Bloom with optimal $k$)
    \item \textbf{Cuckoo Filter}: \texttt{cuckoofilter} v0.5.0 (4-way set-associative)
    \item \textbf{Xor Filter (Xor8)}: \texttt{xorf} v0.11.0 (static/immutable, 8-bit fingerprints)
    \item \textbf{Aleph Filter}: Our implementation (\texttt{aleph-filter} v0.1.0, xxHash3)
  \end{itemize}

  \begin{warningbox}[Xor8 Caveat]
    Xor8 is an \emph{immutable} filter: it is constructed once from a complete key set
    and does not support insertion, deletion, or expansion. Its throughput numbers
    represent the best achievable read performance for a static dataset.
  \end{warningbox}


  % =========================================================================
  \section{Feature Comparison}
  % =========================================================================

  \begin{table}[H]
    \centering
    \caption{Feature comparison across filter types}
    \label{tab:comparison}
    \begin{tabular}{l c c c c c}
      \hline
      \textbf{Feature} & \textbf{Bloom} & \textbf{Cuckoo} & \textbf{Xor8} & \textbf{Aleph} \\
      \hline
      Insert     & \checkmark      & \checkmark       & \texttimes$^{\dagger}$ & \checkmark \\
      Query      & $\bigO{k}$      & $\bigO{1}$       & $\bigO{1}$             & $\bigO{1}$* \\
      Delete     & \texttimes      & \checkmark       & \texttimes             & \checkmark \\
      Expand     & \texttimes      & \texttimes       & \texttimes             & \checkmark \\
      Stable FPR & \texttimes      & \texttimes       & \checkmark$^{\ddagger}$ & \checkmark \\
      \hline
    \end{tabular}

    \smallskip
    {\small * = amortized \quad $\dagger$ = batch build only \quad $\ddagger$ = static set}
  \end{table}


  % =========================================================================
  \section{Throughput Benchmarks}
  % =========================================================================

  All throughput numbers below are from Criterion (100 samples, statistical confidence
  intervals). The custom runner produced consistent results within $\pm$15\%.

  % --- INSERT ---
  \subsection{Insertion Throughput}

  \begin{table}[H]
    \centering
    \caption{Insertion throughput (ns/op, lower is better). Xor8 measures batch \emph{build} time.}
    \label{tab:insert}
    \begin{tabular}{l r r r r}
      \hline
      \textbf{$N$} & \textbf{Aleph} & \textbf{Bloom} & \textbf{Cuckoo} & \textbf{Xor8$^*$} \\
      \hline
      1{,}000   & \textbf{6.3}  & 19.2 & 191.7 & 7.7  \\
      10{,}000  & \textbf{71.9} & 193.7 & 303.7 & 85.5 \\
      100{,}000 & \textbf{1{,}511} & 1{,}974 & 3{,}814 & 1{,}485 \\
      \hline
    \end{tabular}

    \smallskip
    {\small $^*$Xor8 build (batch construction). Per-element amortized.}
  \end{table}

  \begin{insightbox}[Insert Performance]
    The Aleph Filter achieves 66--159 Melem/s insertion throughput, consistently
    outperforming Bloom ($\sim$50 Melem/s) and Cuckoo ($\sim$5--33 Melem/s) due
    to its single-hash design with xxHash3. It is competitive with Xor8's batch
    build at scale.
  \end{insightbox}

  % --- POSITIVE LOOKUP ---
  \subsection{Positive Lookup Throughput}

  Queries for keys that are known to be in the filter (should all return \texttt{true}).

  \begin{table}[H]
    \centering
    \caption{Positive lookup throughput (ns/op, all keys present)}
    \label{tab:lookup-pos}
    \begin{tabular}{l r r r r}
      \hline
      \textbf{$N$} & \textbf{Aleph} & \textbf{Bloom} & \textbf{Cuckoo} & \textbf{Xor8} \\
      \hline
      1{,}000   & 4.5    & 18.4  & 18.3  & \textbf{1.7}  \\
      10{,}000  & 65.2   & 199.3 & 280.4 & \textbf{18.1} \\
      100{,}000 & 3{,}228 & 2{,}050 & 3{,}441 & \textbf{192}  \\
      \hline
    \end{tabular}
  \end{table}

  % --- NEGATIVE LOOKUP ---
  \subsection{Negative Lookup Throughput}

  Queries for keys that are \emph{not} in the filter. This measures the common-case
  query path (early rejection).

  \begin{table}[H]
    \centering
    \caption{Negative lookup throughput (ns/op, no keys present)}
    \label{tab:lookup-neg}
    \begin{tabular}{l r r r r}
      \hline
      \textbf{$N$} & \textbf{Aleph} & \textbf{Bloom} & \textbf{Cuckoo} & \textbf{Xor8} \\
      \hline
      1{,}000   & \textbf{3.8}   & 10.3  & 16.9  & 1.7  \\
      10{,}000  & \textbf{44.8}  & 150.4 & 170.6 & 17.0 \\
      100{,}000 & 2{,}100 & 2{,}673 & 1{,}763 & \textbf{184} \\
      \hline
    \end{tabular}
  \end{table}

  \begin{insightbox}[Lookup Analysis]
    For small to medium datasets ($N \leq 10{,}000$), the Aleph Filter's single-hash
    lookup achieves 222--265 Melem/s on negative queries, outperforming Bloom (66--98 Melem/s)
    and Cuckoo (58--59 Melem/s). Xor8 dominates at $\sim$575 Melem/s due to its cache-friendly
    immutable structure, but cannot support dynamic operations.

    At $N = 100{,}000$, the Aleph Filter's quotient-filter structure (linear probing clusters) shows
    some cache pressure, where Bloom's random-access pattern becomes relatively more efficient.
    This is expected from quotient-filter architectures and aligns with the original paper's observations.
  \end{insightbox}


  % --- DELETE ---
  \subsection{Deletion Throughput}

  Only Aleph and Cuckoo support deletion. Bloom and Xor8 are excluded.

  \begin{table}[H]
    \centering
    \caption{Deletion throughput (ns/op). Bloom and Xor8 do not support deletion.}
    \label{tab:delete}
    \begin{tabular}{l r r}
      \hline
      \textbf{$N$} & \textbf{Aleph} & \textbf{Cuckoo} \\
      \hline
      1{,}000   & \textbf{9.1}  & 33.5  \\
      10{,}000  & \textbf{101.6} & 287.6 \\
      \hline
    \end{tabular}
  \end{table}

  \begin{insightbox}[Deletion]
    Aleph's cluster-aware deletion is $\sim$3$\times$ faster than Cuckoo's
    hash-based eviction at $N = 10{,}000$. This is due to Aleph's compact run-based
    storage: deleting an entry requires only a local shift within the run and cluster,
    whereas Cuckoo must rehash and relocate entries between two candidate buckets.
  \end{insightbox}


  % =========================================================================
  \section{False Positive Rate}
  % =========================================================================

  All filters were configured for a target FPR of 1\%. We measured actual FPR by
  querying 100{,}000 keys not present in the filter.

  \begin{table}[H]
    \centering
    \caption{Empirical false positive rate (\%, target = 1.00\%)}
    \label{tab:fpr}
    \begin{tabular}{l r r r r}
      \hline
      \textbf{$N$} & \textbf{Aleph} & \textbf{Bloom} & \textbf{Cuckoo} & \textbf{Xor8} \\
      \hline
      1{,}000   & \textbf{0.39} & 1.04 & 2.98 & 0.42 \\
      10{,}000  & \textbf{0.43} & 0.99 & 1.98 & 0.39 \\
      100{,}000 & \textbf{0.62} & 1.00 & 2.47 & 0.41 \\
      \hline
    \end{tabular}
  \end{table}

  \begin{insightbox}[FPR Analysis]
    The Aleph Filter consistently achieves a FPR \emph{well below} the 1\% target
    (0.39--0.62\%), demonstrating that the fingerprint-sacrificing expansion mechanism
    does not meaningfully degrade accuracy for these workload sizes. Xor8 achieves
    similarly low rates ($\sim$0.4\%) due to its 8-bit fingerprint design
    ($2^{-8} \approx 0.39\%$).

    The Cuckoo Filter consistently exceeds the target (1.98--2.98\%), likely due to
    the \texttt{cuckoofilter} crate's fixed 12-bit fingerprint width which does not
    precisely match the optimal configuration for a 1\% target.

    Bloom is the most precise relative to its target, landing within $\pm 0.04\%$
    of 1\%, as expected from its well-understood mathematical foundations.
  \end{insightbox}


  % =========================================================================
  \section{Expansion Performance}
  % =========================================================================

  One of the Aleph Filter's unique capabilities is automatic expansion: when the
  load factor exceeds 80\%, the filter doubles its slot count, increments the
  quotient width, and sacrifices one fingerprint bit per entry.

  \begin{table}[H]
    \centering
    \caption{Expansion benchmark: 50{,}000 insertions starting from 64 slots}
    \label{tab:expansion}
    \begin{tabular}{l r}
      \hline
      \textbf{Metric} & \textbf{Value} \\
      \hline
      Amortized insert (with expansion) & 195 ns/op \\
      Total expansions triggered        & 11 \\
      Final capacity                    & 262{,}144 slots \\
      Effective throughput              & 5.1 Melem/s \\
      \hline
    \end{tabular}
  \end{table}

  \begin{warningbox}[Expansion Cost]
    Starting from a 64-slot filter, inserting with expansion (195 ns/op amortized)
    is approximately 13$\times$ slower than pre-sized insertion (15 ns/op). This overhead comes from
    the $\bigO{n}$ re-insertion during each expansion. Users who know the approximate
    dataset size should pre-size the filter via \texttt{AlephFilter::new(expected\_items, fpr)}
    to avoid this cost. However, the filter's ability to expand indefinitely is unique
    among the filters tested and is critical for streaming/unbounded workloads.
  \end{warningbox}


  % =========================================================================
  \section{Conclusion}
  % =========================================================================

  Our Rust implementation of the Aleph Filter demonstrates that the design from the 
  VLDB 2024 paper~\cite{aleph} translates into competitive real-world performance.
  Key findings:

  \begin{enumerate}
    \item \textbf{Fastest insertion} among dynamic filters: 1.3--3$\times$ faster than
          Bloom and 4--30$\times$ faster than Cuckoo, thanks to a single xxHash3 call
          and quotient-filter locality.

    \item \textbf{Strongest negative lookup} at small-to-medium scale: Aleph's early
          rejection via the occupied-flag check achieves 222--265 Melem/s for 
          $N \leq 10{,}000$, beating both Bloom and Cuckoo.

    \item \textbf{Best FPR accuracy}: Measured FPR of 0.39--0.62\% against a 1\% target,
          well within budget and beating Cuckoo's 2--3\% overshoot.

    \item \textbf{Fastest deletion}: $\sim$3$\times$ faster than Cuckoo, the only other
          dynamic filter supporting deletion.

    \item \textbf{Unique expansion capability}: The Aleph Filter is the only filter in this
          comparison that supports automatic, unbounded expansion with fingerprint-bit
          sacrifice, making it the only viable option for streaming workloads where the
          dataset size is not known in advance.
  \end{enumerate}

  The primary tradeoff is at large scale ($N = 100{,}000$), where the quotient-filter
  architecture's linear probing introduces cache pressure during lookups. For
  cache-sensitive, read-heavy workloads with a static dataset, Xor8 remains the optimal
  choice at $\sim$575 Melem/s. For all other use cases, especially those requiring
  dynamic insertion, deletion, and expansion, the Aleph Filter offers the best
  combination of throughput, accuracy, and flexibility.


  \section{Future Work}

  \begin{itemize}
    \item \textbf{SIMD acceleration}: Vectorize run scanning and cluster compaction
          using ARM NEON or x86 AVX2 intrinsics for improved cache utilization at
          large scale.

    \item \textbf{Concurrency}: Add lock-free or fine-grained concurrent access
          for multi-threaded workloads using atomic operations on slot and metadata
          arrays.

    \item \textbf{Compact slot representation}: Reduce the current 72-bit per-slot
          overhead (64-bit slot + 8-bit metadata) by packing metadata flags into
          unused slot bits, potentially halving memory usage.

    \item \textbf{Persistent storage}: Implement memory-mapped I/O for filters
          that exceed RAM, enabling disk-backed probabilistic indexing for
          database systems.

    \item \textbf{Variable fingerprint width}: Allow per-entry fingerprint widths
          instead of the current uniform sacrifice, preserving accuracy for
          frequently-queried entries.
  \end{itemize}

% ============================================================================
% REFERENCES
% ============================================================================
% Uncomment when using biblatex:
% \printbibliography[heading=bibintro, title={References}]

% Manual references for now:
\begin{thebibliography}{9}

\bibitem{carter}
  J.~L.~Carter and M.~N.~Wegman,
  ``Universal Classes of Hash Functions,''
  \textit{Journal of Computer and System Sciences}, vol.~18, no.~2, 1979.
  \url{https://doi.org/10.1016/0022-0000(79)90044-8}

\bibitem{bloom}
  B.~H.~Bloom,
  ``Space/Time Trade-offs in Hash Coding with Allowable Errors,''
  \textit{Communications of the ACM}, vol.~13, no.~7, 1970.
  \url{https://dl.acm.org/doi/10.1145/362686.362692}

\bibitem{chrome-bloom}
  S.~Hess,
  ``Transition Safe Browsing from Bloom Filter to Prefix Set,''
  Chromium Code Review 10896048, September 2012.
  \url{https://codereview.chromium.org/10896048}

\bibitem{mitzenmacher}
  M.~Mitzenmacher and E.~Upfal,
  \textit{Probability and Computing: Randomization and Probabilistic 
  Techniques in Algorithms and Data Analysis},
  Cambridge University Press, 2nd edition, 2017.
  \url{https://www.cambridge.org/core/books/probability-and-computing/7C8B58A4A74AE1A695E0C5F72DB123C8}
  
\bibitem{aleph}
  N.~Dayan, I.-O.~Bercea, and R.~Pagh,
  ``Aleph Filter: To Infinity in Constant Time,''
  \textit{Proceedings of the VLDB Endowment}, vol.~17, no.~11, 2024.
  \url{https://doi.org/10.14778/3681954.3682027}

\bibitem{infinifilter}
  N.~Dayan et al.,
  ``InfiniFilter: Expanding Filters to Infinity and Beyond,''
  \textit{Proceedings of the ACM on Management of Data}, vol.~1, no.~2, 2023.
  \url{https://doi.org/10.1145/3589285}
  
\bibitem{cuckoo}
  B.~Fan, D.~G.~Andersen, M.~Kaminsky, and M.~D.~Mitzenmacher,
  ``Cuckoo Filter: Practically Better Than Bloom,''
  \textit{Proceedings of the 10th ACM International Conference on Emerging 
  Networking Experiments and Technologies (CoNEXT)}, 2014.
  \url{https://doi.org/10.1145/2674005.2674994}
  \url{https://www.cs.cmu.edu/~dga/papers/cuckoo-conext2014.pdf}

\bibitem{quotient}
  M.~A.~Bender et al.,
  ``Don't Thrash: How to Cache Your Hash on Flash,''
  \textit{Proceedings of the VLDB Endowment}, vol.~5, no.~11, 2012.
  \url{https://doi.org/10.14778/2350229.2350275}
\end{thebibliography}
\end{document}
